% ==============================================================
%  Semantic Relaxation Networks:
%  Language Generation as Constraint Stabilization
%  in Admissibility Space
%  Author: Flyxion
% ==============================================================

\documentclass[11pt, a4paper]{article}

% --- Fonts and encoding ---
\usepackage{fontspec}
\setmainfont{TeX Gyre Pagella}
\setsansfont{TeX Gyre Heros}
\usepackage{unicode-math}
\setmathfont{Latin Modern Math}

% --- Microtypography ---
\usepackage{microtype}

% --- Geometry ---
\usepackage[
  top=24mm, bottom=28mm,
  left=30mm, right=30mm,
  headheight=15pt
]{geometry}

% --- Spacing ---
\usepackage{setspace}
\setstretch{1.35}

% --- Paragraph style ---
\usepackage{parskip}
\setlength{\parskip}{0.55em}

% --- Mathematics ---
\usepackage{amsmath}
\usepackage{amsthm}
\usepackage{mathtools}

% --- Theorem environments ---
\newtheoremstyle{srn}
  {1.1em}{0.9em}
  {\normalfont}{}
  {\bfseries\sffamily}{.}
  { }{}

\theoremstyle{srn}
\newtheorem{definition}{Definition}[section]
\newtheorem{principle}{Principle}[section]
\newtheorem{proposition}{Proposition}[section]

\newtheoremstyle{srnremark}
  {0.9em}{0.7em}
  {\normalfont\itshape}{}
  {\itshape\sffamily}{.}
  { }{}

\theoremstyle{srnremark}
\newtheorem{interpretation}{Interpretation}[section]
\newtheorem{remark}{Remark}[section]

% --- Section heading style ---
\usepackage{titlesec}
\titleformat{\section}
  {\large\bfseries\sffamily}
  {\thesection.}{0.7em}{}
\titleformat{\subsection}
  {\normalsize\bfseries\sffamily}
  {\thesubsection.}{0.6em}{}
\titlespacing*{\section}{0pt}{2em}{0.8em}
\titlespacing*{\subsection}{0pt}{1.4em}{0.5em}

% --- Header/footer ---
\usepackage{fancyhdr}
\pagestyle{fancy}
\fancyhf{}
\fancyfoot[C]{\small\thepage}
\renewcommand{\headrulewidth}{0pt}

% --- Hyperlinks ---
\usepackage[numbers]{natbib}
\usepackage[
  colorlinks=true,
  linkcolor=black,
  citecolor=black,
  urlcolor=black
]{hyperref}

% --- Notation macros ---
\newcommand{\Aspace}{\mathcal{A}}           % admissibility space
\newcommand{\Cstruct}{\mathcal{C}}           % constraint/topology structure
\newcommand{\Nhd}{\mathcal{N}}           % semantic neighborhood
\newcommand{\Rop}{\mathcal{R}}           % relaxation operator
\newcommand{\Gop}{\Gamma}              % topology reconfiguration operator
\newcommand{\Phiten}{\varPhi}            % global tension functional
             % instability measure (local)
\newcommand{\dsem}{d_{\mathrm{sem}}}  % semantic distance



% ==============================================================
\begin{document}
% ==============================================================

% --- Title block ---
\begin{center}
  {\LARGE\bfseries\sffamily
    Semantic Relaxation Networks}\\[0.25em]
  {\large\sffamily
    Language Generation as Constraint Stabilization\\
    in Admissibility Space}\\[0.9em]
  {\normalsize\sffamily Flyxion}\\[0.2em]
  {\small\sffamily Independent Researcher}\\[0.7em]
  {\small\sffamily\itshape 2026}
\end{center}

\vspace{0.7em}

% ==============================================================
\begin{abstract}
\noindent
The development of language models over the past decade traces a
compression sequence that has come to feel inevitable: language reduced
to sequence, sequence reduced to tokens, tokens reduced to conditional
probability distributions, and conditional probability treated as a
proxy for cognition itself.
This compression was not an arbitrary choice.
It was a principled engineering projection, one that captured genuine
statistical structure in linguistic data and enabled an extraordinary
range of applications.
But projection always involves loss, and the loss here has been
systematic: what autoregressive factorization discards is precisely
the distributed, recursive, non-sequential character of semantic
structure.

We argue that the sequential decomposition
$p(x) = \prod_i p(x_i \mid x_{<i})$
is computationally convenient rather than cognitively fundamental.
Current language models---including architectures that introduce local
parallelism through diffusion or adaptive tokenization through
byte-level patching---preserve the assumption that text is a
sequential object recovered left-to-right through iterative
approximation.
This assumption, we contend, is not a feature of language as a
cognitive phenomenon.
It is an artefact of treating communicative traces as cognitive
substrate.

We propose Semantic Relaxation Networks (SRNs), a theoretical framework
in which language generation emerges from recursive stabilization over
a partially ordered admissibility manifold rather than from conditional
next-token prediction.
In SRNs, a global tension functional $\Phiten(X)$ replaces the
likelihood objective, semantic configurations evolve through recursive
relaxation operators $\Rop$ driven by constraint gradients, and the
interaction topology $\Cstruct_t$ itself reorganizes dynamically in response
to evolving admissibility structure.
Tokens are not primitive units consumed and emitted by the model.
They are \emph{emergent crystallizations}: local regions of the
admissibility manifold that have stabilized sufficiently to collapse
into communicable symbolic form.

Sequential language, on this account, is not the substrate of
cognition.
It is the observable residue of deeper stabilization dynamics
operating over high-dimensional semantic constraint fields.
\end{abstract}

\newpage

% ==============================================================
\tableofcontents
\vspace{2em}

% ==============================================================
\section{Introduction: Sequentiality as Compression Assumption}
\label{sec:intro}

% --- §1.1 ---
\subsection{The Compression Chain}

There is a genealogy of reductions underlying contemporary language
modelling that is rarely made fully explicit.
It runs as follows.

\emph{Language}, in its cognitive and communicative sense, is a
high-dimensional structured activity involving distributed semantic
commitments, inferential dependencies, prosodic contour, pragmatic
context, and social embedding.
This activity is fundamentally non-linear: meaning is assembled,
revised, anticipated, and recursively stabilized across timescales and
registers that do not reduce to left-to-right symbol emission.

The first compression identifies language with \emph{sequence}.
Written text is sequential; the temporal unfolding of speech is
sequential; it is therefore tempting---and technically tractable---to
model language as a sequence of discrete symbols.
This move is ancient in formal linguistics and natural language
processing alike.
Its power is undeniable.
Its cost is that the relational, distributed, and recursive structure
of semantic coherence must be reconstructed from positional proximity
alone.

The second compression identifies sequence with \emph{tokens}.
Sequences must be discretized to be processed by finite machines.
Whether the unit of discretization is a word, a subword piece, a byte,
or a patch of bytes, the operational assumption is that meaning can be
adequately represented by a finite vocabulary of symbols arrayed in
linear order.

The third compression identifies tokens with \emph{conditional
probabilities}.
The chain rule of probability yields
\[
  p(x_1, x_2, \ldots, x_n)
  \;=\; \prod_{i=1}^n p\!\left(x_i \,\middle|\, x_{<i}\right),
\]
an exact factorization of any joint distribution over sequences.
The autoregressive language model implements this factorization as its
generative mechanism: at each position, predict the next token given
all previous tokens.
This is not merely a modelling convention.
It becomes an architectural commitment, shaping every aspect of
training, inference, and evaluation.

The fourth and most consequential compression identifies conditional
prediction with \emph{cognition itself}.
If a system that predicts the next token with sufficient accuracy also
translates, reasons, writes code, and answers questions, the temptation
arises to conclude that intelligence simply is high-quality conditional
prediction.
This inference is not obviously wrong.
But it conflates the communicative trace of cognition---sequential
language---with the underlying cognitive processes that generate and
comprehend it.

\subsection{Partial Destabilizations}

Contemporary architectures have already begun, implicitly, to
destabilize the most rigid layers of this compression.

Denoising diffusion models applied to language \citep{ho2020ddpm,
austin2021discrete, lou2024diffusion} replace left-to-right generation
with iterative denoising: a corrupted sequence is incrementally
restored toward a coherent target.
Diffusion weakens strict causal ordering and introduces a form of
parallel structure.
But it preserves the fundamental assumption that generation is
\emph{sequence recovery}: there is always an underlying ordered text
toward which the model converges.
The topology of interaction---which positions attend to which---remains
essentially sequence-indexed.

Block diffusion \citep{arriola2025block} interpolates between
autoregressive and diffusion regimes, allowing parallel generation
within local windows while preserving sequential ordering between
windows.
This is a meaningful relaxation, but sequentiality is merely
decomposed rather than questioned.

The Byte Latent Transformer \citep{pagnoni2025blt} and its faster
variant \citep{kallini2026fastblt} challenge the tokenization layer
by operating at byte level with dynamic entropy-based patching.
High-entropy byte sequences receive finer-grained attention; redundant
regions are compressed into larger patches.
BLT thereby weakens the fixed-vocabulary assumption and introduces
content-adaptive processing granularity.
Yet the architecture still generates text as a sequential output,
proceeding through the sequence with a modified notion of step size.

These developments are not failures.
They represent genuine conceptual progress: each one loosens a
different constraint in the compression chain.
Diffusion loosens strict ordering.
Block diffusion loosens strict left-to-right causal dependency.
BLT loosens fixed tokenization.

But none of them questions whether \emph{sequentiality itself} is the
right ontology for language generation.
They are, in the terminology we will develop, \emph{transitional
architectures}: they operate within admissibility geometry without
yet recognizing that geometry as their proper domain.

\subsection{Thesis}

This paper proposes a more fundamental reorientation.

We argue that sequential text is not the primitive substrate of
language generation but its \emph{observational residue}: a
low-dimensional projection of higher-dimensional stabilization
processes operating over semantic constraint fields.
A language model that generates text by predicting the next token is
performing an operation analogous to reading a shadow and inferring the
object that cast it---accurate within limits, but ontologically
inverted.

The framework we develop, Semantic Relaxation Networks, replaces the
next-token prediction objective with a \emph{recursive stabilization
dynamics} over a partially ordered admissibility manifold.
Generation is not emission.
It is the collapse of locally stabilized regions of semantic
configuration space into communicable symbolic form.

The central distinction organizing the paper is this:

\begin{center}
\begin{tabular}{ll}
\textbf{Transformers predict.} & \textbf{SRNs stabilize.}\\
Prediction is local and sequential. & Stabilization is distributed and recursive.\\
The unit is the token. & The unit is the admissibility region.\\
Topology is fixed by sequence. & Topology reorganizes under constraint.\\
Hallucination is misassignment. & Hallucination is premature collapse.
\end{tabular}
\end{center}

\noindent The paper proceeds as follows.
Section~\ref{sec:failure} identifies the structural failure modes
of sequential ontology more precisely.
Section~\ref{sec:admissibility} defines admissibility spaces and
their formal properties.
Section~\ref{sec:dynamics} develops the recursive relaxation dynamics
and companion equations.
Section~\ref{sec:topology} introduces dynamic semantic topology
formation, the key architectural departure from transformers.
Section~\ref{sec:verification} treats stabilization, verification
operators, and the reinterpretation of hallucination.
Section~\ref{sec:emergent} presents the inversion: emergent
sequentiality as projection residue.
Section~\ref{sec:rsvp} situates SRNs within the RSVP, TARTAN, and
CLIO frameworks.
Section~\ref{sec:implications} draws implications for AI, cognitive
science, and physics.
Section~\ref{sec:conclusion} concludes.


\subsection{The SRN Object Hierarchy}
\label{subsec:hierarchy}

Before the mathematics becomes dense, we introduce the five primary
formal objects of the SRN framework and fix their roles explicitly.
The single largest source of conceptual ambiguity in constraint-based
generation proposals is the conflation of \emph{dynamical} objects
(those that evolve during relaxation) with \emph{residual} objects
(those that persist after stabilization is complete).
The following definition and table separate these cleanly.

\begin{definition}[SRN Object Hierarchy]
\label{def:hierarchy}
The SRN framework is organized around five primary formal objects:
\begin{enumerate}
  \item[(O1)] The \emph{semantic configuration} $X_t \in \mathcal{X}$:
    the current state of the partially stabilized semantic field at
    relaxation step $t$.
  \item[(O2)] The \emph{interaction topology} $\Cstruct_t$: the
    weighted directed graph encoding current coupling strengths between
    semantic regions, evolving under the topology operator $\Gop$.
  \item[(O3)] The \emph{stabilized dependency structure} $G(X^*)$: the
    directed acyclic graph of communicative precedence constraints
    residual in the fully stabilized configuration $X^*$; a static
    object extracted once stabilization is complete.
  \item[(O4)] The \emph{verification hierarchy} $\mathcal{V} = \{V_k\}$:
    the multi-scale family of admissibility checking operators that
    validate local stabilizations against global constraints and
    activate rollback when violations are detected.
  \item[(O5)] The \emph{communicative pressure field} $\mathcal{P}$:
    the weighted preference function over ordered region pairs encoding
    grammatical, information-structural, and prosodic linearization
    conventions; applied during projection, not during relaxation.
\end{enumerate}
\end{definition}

\begin{center}
\renewcommand{\arraystretch}{1.3}
\begin{tabular}{lllll}
\hline
\textbf{Object} & \textbf{Notation} & \textbf{Role} &
\textbf{Type} & \textbf{Section} \\
\hline
Semantic configuration & $X_t$ &
  Current field state & Dynamical & \S\ref{sec:admissibility} \\
Interaction topology & $\Cstruct_t$ &
  Coupling graph & Dynamical & \S\ref{sec:topology} \\
Stabilized DAG & $G(X^*)$ &
  Dependency skeleton & Residual & \S\ref{sec:emergent} \\
Verification hierarchy & $\mathcal{V}$ &
  Admissibility checking & Dynamical & \S\ref{sec:verification} \\
Pressure field & $\mathcal{P}$ &
  Linearization preferences & Static & \S\ref{sec:projection} \\
\hline
\end{tabular}
\end{center}

\medskip
\noindent
The key distinction is between $\Cstruct_t$ and $G(X^*)$.
The interaction topology $\Cstruct_t$ reorganizes continuously during
relaxation under the operator $\Gop$ in response to constraint
gradients; it is a computational scaffold, not a semantic claim.
The stabilized DAG $G(X^*)$ is extracted once from the converged
configuration and encodes the genuine semantic dependency structure of
the output; it is the object over which projection operates.
Conflating these two objects---treating the relaxation graph as
identical to the dependency structure---is the source of most
misreadings of constraint-based generation proposals.

% ==============================================================
\section{The Failure of Sequential Ontology}
\label{sec:failure}

\subsection{What Autoregressive Factorization Discards}

The autoregressive factorization
\[
  p(x) \;=\; \prod_{i=1}^n p(x_i \mid x_{<i})
\]
is mathematically exact as a decomposition of any joint distribution
over sequences.
This is a theorem, not an assumption.
But the \emph{implementation} of this factorization as a generative
mechanism involves choices that are not forced by the mathematics.

The most consequential such choice is the \emph{Markovian
approximation}: the joint context $x_{<i}$ is compressed into a
finite-dimensional representation by the model.
For transformer architectures \citep{vaswani2017attention}, this
representation is the softmax-weighted combination of previous-position
key-value pairs.
The attention mechanism is extraordinarily expressive, but it operates
over a fixed interaction structure determined by sequence position.
Two tokens interact in proportion to their attention weight, which is a
function of their content and their relative position---not of their
semantic compatibility or constraint relationship.

The result is a subtle but structurally important limitation.
A transformer's interaction graph is \emph{position-indexed}.
Semantic coherence must be inferred \emph{through} positional
relationships rather than being represented \emph{directly}.
Long-range semantic dependencies are, in principle, representable
through attention; in practice they compete with local positional
patterns that dominate the attention distribution.

More fundamentally, the autoregressive framework forces a
\emph{strict temporal ordering} on the generative process.
Each token is produced after all preceding tokens are fixed.
This means the model cannot revise a token once emitted, cannot
simultaneously stabilize distant regions of the output, and cannot
represent genuine semantic feedback between non-adjacent regions during
generation.

These are not merely engineering limitations to be overcome with
better hardware.
They are \emph{ontological} commitments built into the factorization
itself.

\subsection{Hallucination as Structural Symptom}

The phenomenon of hallucination---the confident generation of
semantically plausible but factually or inferentially incorrect
text---is typically treated as a failure mode to be mitigated through
better training, retrieval augmentation, or output verification.

Under sequential ontology, hallucination has a straightforward
description: the model assigned high conditional probability to an
incorrect continuation.
The remedy is therefore to improve the calibration of the conditional
distribution.

We propose that this description, while locally accurate, misidentifies
the structural source of the failure.

Hallucination is not primarily a calibration problem.
It is a \emph{topological instability event}: a premature local
stabilization that fails to maintain global admissibility.

The autoregressive model commits to each token irrevocably.
Once a false commitment is made, subsequent tokens are conditioned on
a flawed context, and the model's only recourse is to continue
generating text that is locally coherent with the erroneous token
rather than globally coherent with the semantic task.
Sequential generation has no natural mechanism for detecting that a
locally plausible choice has produced a globally inadmissible
configuration.

This is not a failure of attention or memory.
It is a consequence of treating generation as irreversible sequential
emission rather than as recursive stabilization with rollback capacity.

\subsection{Cognition Is Not Sequential}

Perhaps the deepest problem with sequential ontology is the implicit
claim it makes about cognition.

Human language comprehension and production are not sequential
processes at the level of cognitive implementation.
Decades of psycholinguistic research establish that comprehension
involves massive parallelism, predictive processing, and global
constraint satisfaction \citep{pearl1988probabilistic}.
Readers do not process text token-by-token, assigning conditional
probabilities to each word given its predecessors.
They construct meaning through distributed inference across multiple
linguistic levels simultaneously, frequently revising interpretations
in light of later information.

Production, similarly, involves hierarchical planning across timescales
that do not reduce to incremental token emission
\citep{smolensky1986harmony, hinton1983optimal}.
The sequential character of speech and writing is a \emph{interface
constraint}---the communication channel is one-dimensional---not a
property of the underlying generative process.

Free energy and predictive processing frameworks
\citep{friston2010free, friston2006free} have proposed that neural
computation is fundamentally a process of minimizing prediction error
(or equivalently, minimizing variational free energy) across hierarchical
generative models.
These frameworks are explicitly non-sequential: the brain simultaneously
maintains predictions at multiple levels of abstraction and updates them
in parallel.

LeCun's world model framework \citep{lecun2022path} similarly
argues that intelligence requires structured world models that support
explicit planning and constraint satisfaction rather than purely
reactive next-step prediction.

Autoregressive language models, for all their empirical power, adopt
an architecture that is in tension with these converging insights from
cognitive science and neuroscience.
They are, in a precise sense, \emph{communication channel models}
rather than \emph{cognitive substrate models}.

The following sections develop an alternative.

% ==============================================================
\section{Admissibility Spaces and Semantic Persistence}
\label{sec:admissibility}

\subsection{Beyond Probability}

The standard framework for language generation measures the quality of
a sequence by its probability under a learned distribution.
Generation maximizes (or samples from) $p(x)$.
Evaluation metrics are cast in terms of perplexity, likelihood, or
functions thereof.

We propose to replace probability as the primary object with a
different notion: \emph{admissibility}.

Admissibility is not equivalent to likelihood, though the two are
related.
A highly probable sequence may be inadmissible: it may be locally
fluent but globally incoherent, or fluent under current conditions but
unable to survive recursive extension.
Conversely, an admissible configuration may not be the most probable
continuation at any given position; it may be a semantically
structured but stylistically unusual configuration that maintains
coherence under perturbation and extension.

The distinction is essential.
Probability models ask: \emph{does this continuation fit locally?}
Admissibility asks: \emph{can this configuration survive recursive
future stabilization?}

This is a fundamentally different question, and answering it requires
a fundamentally different formal apparatus.


\subsection{Semantic Regions as Primitive Objects}

Before defining admissibility space, we must specify what its elements
are.
The framework has so far referred freely to ``sub-configurations,''
``regions,'' and ``patches'' of the semantic configuration $X$.
These terms require a formal placeholder definition, since the whole
apparatus of tension functionals, topology operators, and projection
maps is defined over them.

\begin{definition}[Semantic Region]
\label{def:region}
A \emph{semantic region} $X_i$ is a locally coherent subconfiguration
of the semantic configuration space $\mathcal{X}$ satisfying:
\begin{enumerate}
  \item[(i)] \textbf{Internal cohesion.}
    The internal constraint coupling of $X_i$---the mean edge weight
    among nodes within $X_i$ under the current topology
    $\Cstruct_t$---exceeds its external coupling to all other regions
    over some relaxation interval $[t_0, t_0 + \delta]$.
    Formally, $\bar{\omega}^{\mathrm{int}}_i(t) > \bar{\omega}^{\mathrm{ext}}_i(t)$
    for $t \in [t_0, t_0 + \delta]$.
  \item[(ii)] \textbf{Semantic coherence.}
    The local instability measure $S_i(X) < \tau_{\mathrm{region}}$
    for a region-formation threshold $\tau_{\mathrm{region}} \leq \tau$,
    meaning the region does not contain unresolved internal contradiction.
  \item[(iii)] \textbf{Temporal persistence.}
    $X_i$ maintains conditions (i) and (ii) across at least one
    relaxation step, distinguishing it from transient fluctuations in
    the constraint graph.
\end{enumerate}
\end{definition}

\begin{remark}
This definition is deliberately operational rather than ontologically
final.
Semantic regions are not fixed objects: they form, merge, split, and
dissolve as the relaxation dynamics evolve and the topology
$\Cstruct_t$ reorganizes.
A region that satisfies (i)--(iii) at step $t$ may fail to satisfy
them at step $t+1$ if a topology reconfiguration increases its
external coupling or introduces internal contradiction.
The definition therefore characterizes regions as \emph{dynamically
stable patterns} in the constraint graph rather than as static
constituents of a fixed parse.

This flexibility is intentional.
Natural semantic units---words, phrases, propositions, discourse
segments---do not correspond to fixed-granularity objects.
They are better understood as dynamically stable constraint clusters
whose boundaries shift with context, attention, and communicative
purpose.
The SRN framework encodes this by treating regions as emergent
features of the relaxation process rather than as primitive inputs.
\end{remark}

\subsection{Formal Definition}

We introduce the central formal object of the paper.

\begin{definition}[Admissibility Space]
\label{def:admissibility}
Let $\mathcal{X}$ be the space of semantic configurations.
An \emph{admissibility space} $\Aspace \subseteq \mathcal{X}$ is a
dynamically evolving, partially ordered persistence structure on
$\mathcal{X}$ such that each element $X \in \Aspace$ satisfies the
following three conditions:

\begin{enumerate}
  \item[\textbf{(C1)}] \textbf{Local Constraint Coherence.}
    For every sub-configuration $X_i \subseteq X$ and its evolving
    semantic neighborhood $\Nhd_i \subseteq \Aspace$, the local constraint
    tension $S_i(X) < \tau$ for some coherence threshold $\tau > 0$.
    That is, $X_i$ does not generate unresolvable local contradiction
    with its immediate semantic context.

  \item[\textbf{(C2)}] \textbf{Recursive Extensibility.}
    For every admissible extension operator $\mathcal{E}$ consistent
    with the global constraint structure $\Cstruct$, there exists an
    admissible extension $X' = \mathcal{E}(X)$ such that $X'
    \in \Aspace$ and the global tension functional $\Phiten(X') \leq \Phiten(X)$.
    That is, $X$ can be continued without necessarily increasing total
    semantic tension.

  \item[\textbf{(C3)}] \textbf{Perturbation Stability under Field
    Deformation.}
    $X$ remains admissible under moderate reorganization of the
    semantic interaction topology $\Cstruct$.
    Formally, for perturbations $\delta\Cstruct$ below a topology deformation
    threshold $\epsilon$, we require $X \in \Aspace(\Cstruct + \delta\Cstruct)$.
    Stability here is with respect to \emph{topological}, not merely
    statistical, perturbation.
\end{enumerate}
\end{definition}

\begin{interpretation}
Condition (C1) is the local analogue of syntactic and semantic
well-formedness: the configuration does not generate immediate
contradiction within its local semantic neighborhood.
Condition (C2) is the more demanding and more philosophically novel
condition: admissibility is not merely a property of the current state
but of its \emph{future persistence}.
An admissible configuration must be one from which further coherent
stabilization is possible.
This separates admissibility from likelihood: a high-probability
continuation may close off subsequent extension.
Condition (C3) ensures that admissibility is robust to semantic
topology reorganization---the model's ability to dynamically
reconfigure interaction structure (Section~\ref{sec:topology}) must
not invalidate previously stabilized configurations.
\end{interpretation}

\begin{remark}
The term ``partially ordered'' is important.
Admissibility space is not a simple inclusion structure.
Configurations are ordered by their degree of stabilization, their
depth of recursive extensibility, and their robustness to perturbation.
This ordering is partial because not all configurations are comparable:
two configurations may each be admissible without either being more
admissible than the other.
The space is therefore a directed partially ordered set (a
\emph{poset}) rather than a total order.
\end{remark}

\subsection{Admissibility and RSVP Accessibility}

For readers familiar with the Relativistic Scalar-Vector Plenum (RSVP)
framework, admissibility space admits a natural interpretation as the
semantic analogue of accessibility structure under entropy-bounded
persistence \citep[cf.]{barbour2003end}.
In RSVP terms, a configuration $X$ is admissible when it lies within
the accessible region of semantic field space consistent with
maintaining low entropic tension under recursive extension.

This connection is developed more fully in Section~\ref{sec:rsvp}.
For the purposes of the present formal development, admissibility
space should be understood on its own terms: as a persistence-theoretic
object defined by the three conditions above, independent of any
commitment to RSVP cosmology.


\subsection{Admissibility, Truth, and External Verification}

A possible misreading of the framework must be addressed directly.
Admissibility, as defined through conditions (C1)--(C3), might appear
to reduce to \emph{self-consistency}: a configuration is admissible if
it does not contradict itself and can be extended without
self-contradiction.
This would make SRNs glorified coherence engines---systems that produce
fluent, internally consistent text without any anchor to truth or
reality.

This reading is incorrect, but refuting it requires making explicit
what the framework implies about false content.

A false configuration---one whose semantic commitments do not
correspond to states of affairs in the world---may occupy a locally
admissible basin.
Internal coherence does not guarantee truth.
A well-constructed fiction, a plausible lie, or a sophisticated
confabulation may satisfy (C1) locally and even satisfy (C2) in
the short term.

However, the \emph{recursive extensibility} condition (C2) is more
demanding than it first appears.
Extensibility is not evaluated in isolation.
It is evaluated under continued interaction with external verification
fields: sensory input, interlocutor responses, documentary evidence,
and logical inference from established commitments.
These external fields function as additional constraint sources that
are continuously integrated into the admissibility computation.

\begin{principle}[Admissibility Under External Pressure]
\label{prin:truth}
Let $\mathcal{F}_{\mathrm{ext}}$ denote the set of external
verification fields (empirical, inferential, social) available during
generation.
A configuration $X$ that is transiently admissible in isolation may
become inadmissible when the global tension functional is augmented
with external constraint terms:
\[
  \Phiten_{\mathrm{ext}}(X)
  \;=\;
  \Phiten(X)
  \;+\;
  \mu \sum_{k} \Psi_k(X, \mathcal{F}_k),
\]
where $\Psi_k(X, \mathcal{F}_k)$ measures the tension between
configuration $X$ and external field $\mathcal{F}_k$, and $\mu > 0$
weights external against internal constraint.
False configurations generate high $\Psi_k$ over time as their
predictions conflict with field observations, progressively increasing
$\Phiten_{\mathrm{ext}}$ and destabilizing the configuration.
\end{principle}

\begin{interpretation}
Principle~\ref{prin:truth} has several important consequences.

First, \emph{hallucination and systematic falsehood are structurally
distinct under SRN}.
Hallucination is premature crystallization---a local collapse before
global admissibility is verified (Principle~\ref{prin:hallucination}).
Systematic falsehood is a configuration that is locally stable in
isolation but progressively destabilized by external fields over time.
Both fail admissibility, but at different timescales and for different
structural reasons.

Second, \emph{truth-tracking is a dynamical property}, not a static
label.
A configuration does not simply ``have'' a truth value; it maintains
or loses admissibility as it interacts with external verification
fields.
This is consistent with the epistemological insight that truth is a
relationship between representations and the world, not an intrinsic
property of the representation alone.

Third, \emph{the SRN framework is not committed to a correspondence
theory of truth as a foundational axiom}.
It requires only that external verification fields exist and exert
constraint pressure on configurations that conflict with them.
What those fields consist of---empirical observation, logical inference,
social consensus---is left to the domain of application.

Fourth, this account provides a principled explanation of why
language models trained on text alone can produce coherent falsehoods.
Without external verification fields integrated into $\Phiten$, the
model minimizes internal tension only.
Retrieval augmentation, tool use, and grounding mechanisms are, in SRN
terms, methods for incorporating $\mathcal{F}_{\mathrm{ext}}$ into
the tension functional during generation.
\end{interpretation}


\subsection{The Role of Probability in the SRN Framework}

The SRN framework does not reject probability theory.
This point deserves explicit statement, because the paper's critique of
autoregressive factorization might be misread as a claim that
conditional probability is incorrect or that likelihood is a wrong
objective.
That is not the position.

The SRN position is more precise: conditional probability is a
\emph{projection statistic over stabilization traces}, not the
primitive ontology of generation.

When a fully stabilized configuration $X^*$ is projected via $\pi$
into a sequential output, that output has a probability distribution
induced by the pressure field $\mathcal{P}$ over the space of linear
extensions $\mathrm{Lin}(G(X^*))$.
That distribution is well-defined and useful for many purposes:
evaluation, comparison, and sampling.
Conditional probability $p(x_i \mid x_{<i})$ remains a
locally valid estimator of how often one region crystallizes before
another in a given context---a useful approximation of pressure-field
dynamics.

What the framework disputes is the claim that this probability is the
\emph{primary object} of generation---that generation consists in
computing $p(x_i \mid x_{<i})$ and sampling from it, rather than in
stabilizing admissibility structure and projecting it.
The distinction is between: optimizing a probability model of
sequential traces, and stabilizing a constraint geometry of which
sequential traces are projections.

Both approaches can produce similar outputs in practice.
The difference emerges at the failure modes: autoregressive
probability models fail at garden-path disambiguation, long-range
admissibility, recursive extensibility, and global consistency
precisely because they optimize a projection statistic rather than
the underlying geometry.
SRNs propose to optimize the geometry directly, treating probability
as one observable shadow of a deeper stabilization process.

% ==============================================================
\section{Recursive Relaxation Dynamics}
\label{sec:dynamics}

\subsection{Tension Functionals}

The operational core of the SRN framework is a family of
\emph{semantic tension measures} that quantify the degree to which a
configuration or pair of configurations fails to satisfy admissibility
conditions.

\subsubsection*{Local Instability}

For a sub-configuration $X_i$ with evolving semantic neighborhood
$\Nhd_i$, define the \emph{local instability measure}:
\[
  S_i(X)
  \;=\;
  H\!\left(X_i \,\middle|\, \Nhd_i\right)
  \;+\;
  \lambda\, D_i(X),
\]
where $H(X_i \mid \Nhd_i)$ denotes the conditional uncertainty of $X_i$
given its semantic neighborhood (the degree to which $X_i$ is
unpredicted by its constraint context), $D_i(X)$ is a
\emph{contradiction penalty} measuring direct semantic incompatibility
between $X_i$ and elements of $\Nhd_i$, and $\lambda > 0$ is a
weighting parameter.

\begin{interpretation}
Local instability is not equivalent to local entropy.
A configuration may be statistically unusual---high conditional
entropy---yet semantically coherent and therefore admissible.
Conversely, a configuration may be locally probable yet semantically
contradictory with its neighborhood.
The contradiction penalty $D_i$ captures this: it registers direct
inferential or semantic incompatibility rather than mere statistical
unpredictability.
\end{interpretation}

\subsubsection*{Semantic Tension Between Regions}

For a pair of sub-configurations $(X_i, X_j)$, define the
\emph{semantic tension}:
\[
  T(i,j)
  \;=\;
  \omega_{ij}\, \dsem(X_i, X_j)
  \;+\;
  \mu\, \kappa_{ij},
\]
where $\omega_{ij} \geq 0$ is a coupling strength reflecting the
degree to which regions $i$ and $j$ are semantically interdependent,
$\dsem(X_i, X_j)$ is a semantic incompatibility distance on
$\mathcal{X}$, and $\kappa_{ij}$ is a \emph{recursive constraint
conflict} term measuring unresolved dependency between the two regions
under future extension.

\subsubsection*{Global Tension Functional}

The global tension functional integrates local instability and pairwise
semantic tension:
\[
  \Phiten(X)
  \;=\;
  \sum_{i} S_i(X)
  \;+\;
  \sum_{i < j} T(i,j).
\]
Generation in the SRN framework consists in minimizing $\Phiten(X)$
over the admissibility space $\Aspace$ subject to the constraint structure
$\Cstruct_t$.


\subsection{Possible Realizations of Semantic Distance and Contradiction Metrics}

The tension functional $\Phiten(X) = \sum_i S_i(X) + \sum_{i<j} T(i,j)$
depends on three quantities that have been left intentionally abstract:
the semantic distance $\dsem(X_i, X_j)$, the contradiction penalty
$D_i(X)$, and the recursive constraint conflict $\kappa_{ij}$.
This abstraction is deliberate: the SRN framework is designed to be
\emph{substrate-agnostic}, specifying the functional role of each
quantity without committing to a particular implementation.

However, a skeptical reader will reasonably ask what category of
mathematical object each quantity actually is.
This subsection provides concrete candidate realizations, organized
by the kind of substrate being modeled.

\subsubsection*{Semantic Distance $\dsem(X_i, X_j)$}

The semantic distance measures incompatibility between regions $X_i$
and $X_j$ in the configuration space $\mathcal{X}$.
Candidate realizations include:

\begin{itemize}
  \item \textbf{Embedding divergence.}
    If semantic regions are represented as distributions over a latent
    space, $\dsem$ can be instantiated as KL divergence, Wasserstein
    distance, or cosine dissimilarity between embedding vectors.
    This connects SRNs directly to existing representation learning
    methods.

  \item \textbf{Graph Laplacian distance.}
    If $\mathcal{X}$ carries a graph structure (e.g., a knowledge graph
    or semantic network), $\dsem$ can be defined via the commute-time
    distance or resistance distance on the graph Laplacian $L$:
    $\dsem(X_i, X_j) = (e_i - e_j)^\top L^+ (e_i - e_j)$, where
    $L^+$ is the Moore-Penrose pseudoinverse.

  \item \textbf{Logical incompatibility.}
    In a formal-language setting, $\dsem(X_i, X_j)$ can be the
    minimum number of axiom applications needed to derive a
    contradiction from the joint assertion of $X_i$ and $X_j$, or
    $\infty$ if they are jointly consistent.

  \item \textbf{Category-theoretic obstruction.}
    If configurations are objects in a category and compatibility is
    modeled by morphisms, $\dsem$ can measure the obstruction to the
    existence of a colimit or pushout for the diagram $(X_i, X_j)$---a
    sheaf-cohomological incompatibility in the sense of
    \v{C}ech cohomology on the covering of semantic space.
\end{itemize}

\subsubsection*{Contradiction Penalty $D_i(X)$}

The contradiction penalty $D_i(X)$ measures direct semantic
incompatibility between region $X_i$ and its neighborhood
$\Nhd_i^t$.
Unlike $\dsem$, which is symmetric and pairwise, $D_i$ is a local
quantity measuring how well $X_i$ fits its current context.
Candidate realizations include:

\begin{itemize}
  \item \textbf{Constraint satisfaction deficit.}
    In a constraint satisfaction problem (CSP) formulation, $D_i$
    counts the number of constraints involving $X_i$ that are violated
    given the current assignments of its neighbors.
    This gives a SAT-style incompatibility count directly analogous to
    the unsatisfied-clause count in MAX-SAT.

  \item \textbf{Predictive error.}
    In a neural implementation, $D_i$ can be instantiated as the
    reconstruction error of $X_i$ under a model trained to predict
    each region from its semantic neighborhood---a local predictive
    coding error.
    This connects $D_i$ directly to the free-energy principle
    \citep{friston2010free}.

  \item \textbf{Graph potential violation.}
    If the interaction graph $\Cstruct_t$ carries potential functions
    on edges (as in a Markov Random Field), $D_i$ is the sum of
    pairwise potentials $\phi(X_i, X_j)$ for all neighbors $j \in
    \Nhd_i^t$ that are in a conflicting state with $X_i$.
\end{itemize}

\subsubsection*{Recursive Constraint Conflict $\kappa_{ij}$}

The recursive constraint conflict $\kappa_{ij}$ captures unresolved
dependency between regions $X_i$ and $X_j$ under future extension---a
forward-looking incompatibility that is not captured by the current
state alone.
This is the most novel of the three quantities and the hardest to
instantiate directly.
Candidate approaches include:

\begin{itemize}
  \item \textbf{Extension divergence.}
    $\kappa_{ij}$ can be defined as the divergence between the
    distribution of admissible extensions reachable from $X_i$ alone
    and those reachable from $X_j$ alone: regions that would ``pull''
    future stabilization in incompatible directions exhibit high
    $\kappa_{ij}$ even if their current states are locally compatible.

  \item \textbf{Rollback correlation.}
    In a trained system with rollback statistics, $\kappa_{ij}$ can be
    estimated empirically as the frequency with which accepting region
    $X_j$ later triggers rollback of region $X_i$ (or vice versa)
    across a corpus of generation traces.
    This makes $\kappa_{ij}$ a learnable quantity from interaction data.

  \item \textbf{Commitment interference.}
    In formal-language settings, $\kappa_{ij}$ measures the reduction
    in the set of admissible completions caused by jointly committing
    to both $X_i$ and $X_j$:
    $\kappa_{ij} = |\mathrm{Ext}(\emptyset)| - |\mathrm{Ext}(\{X_i,
    X_j\})|$, normalized by the baseline extension count.
\end{itemize}

\begin{remark}
The framework is intentionally agnostic among these realizations.
Different substrates will favor different instantiations:
neural text generation favors embedding divergence and predictive
error; formal verification favors logical incompatibility and
constraint satisfaction deficit; biological cognition favors
predictive coding error and rollback correlation.
The SRN framework specifies the \emph{functional roles} of
$\dsem$, $D_i$, and $\kappa_{ij}$ in the dynamics; which
mathematical objects fill those roles is a question for the domain
of application.
\end{remark}

\begin{principle}[Generative Objective]
\label{prin:objective}
Language generation in an SRN is not optimization of token-level
likelihood. It is minimization of semantic tension across the
admissibility manifold:
\[
  X^* \;=\; \arg\min_{X \in \Aspace} \Phiten(X).
\]
Sequential tokens are emitted only when local regions of $X^*$ achieve
sufficient stabilization---that is, when $S_i(X^*) < \tau$ and the
region $X_i$ can no longer benefit from further relaxation.
\end{principle}



\subsection{The Vector-Valued Tension Functional}

The scalar tension functional $\Phiten(X) = \mathbf{w}^\top
\vec{\Phiten}(X)$ used throughout this paper is an operational
projection of a richer underlying structure.
We make this explicit here, for two reasons: first, several sections
already implicitly assume multi-objective tension geometry
(the epistemic decomposition of §\ref{sec:implications}, the
sycophancy failure analysis, the distinction between local and
global admissibility criteria); second, the extension reveals that
scalar optimization is itself a projection operation, mirroring the
paper's broader thesis that sequentiality, tokens, and grammar are
all projections of higher-dimensional admissibility structure.

\begin{definition}[Vector-Valued Tension Functional]
\label{def:vecphi}
The \emph{vector-valued tension functional} is a map
\[
  \vec{\Phiten} : \mathcal{X} \;\longrightarrow\; \mathbb{R}^n,
\]
whose $k$-th component $\Phiten_k(X)$ measures the total tension
along the $k$-th admissibility dimension.
Natural dimension assignments include:
\begin{enumerate}
  \item[(V1)] $\Phiten_{\mathrm{sem}}$: internal semantic
    incompatibility, as defined by the local instability and pairwise
    tension terms in §\ref{sec:dynamics}.
  \item[(V2)] $\Phiten_{\mathrm{ext}}$: external verification
    tension, as defined by the field terms $\Psi_k(X,
    \mathcal{F}_k)$ in Principle~\ref{prin:truth}.
  \item[(V3)] $\Phiten_{\mathrm{proj}}$: projection instability,
    measuring sensitivity of linearization to pressure-field
    perturbation (related to the failure metric $E_{\mathrm{proj}}$).
  \item[(V4)] $\Phiten_{\mathrm{conv}}$: conversational tension,
    as introduced in the epistemic stability analysis
    (§\ref{sec:implications}).
  \item[(V5)] $\Phiten_{\mathrm{adv}}$: adversarial
    under-pressure, measuring insufficiency of constraint application
    (component (E3) of Definition~\ref{def:epistemic}).
\end{enumerate}
The scalar tension functional used throughout the main development
is the weighted projection:
\begin{equation}
  \label{eq:scalar_proj}
  \Phiten(X) \;=\; \mathbf{w}^\top \vec{\Phiten}(X)
  \;=\; \sum_{k=1}^n w_k\, \Phiten_k(X),
\end{equation}
for a task-specific weight vector $\mathbf{w} \in \mathbb{R}^n_{\geq 0}$.
\end{definition}

\begin{interpretation}
Equation~\eqref{eq:scalar_proj} is itself a projection operation.
The scalar tension landscape $\Phiten : \mathcal{X} \to \mathbb{R}$
that the relaxation operator descends is not the primitive object; it
is a one-dimensional shadow of the higher-dimensional admissibility
pressure structure $\vec{\Phiten} : \mathcal{X} \to \mathbb{R}^n$.

This mirrors the paper's central thesis at the level of its own
formalism.
Sequential language is a projection of admissibility geometry;
scalar tension is a projection of multi-dimensional admissibility
pressure.
The same compression operation---discarding dimensional structure in
favour of a single tractable scalar---appears at both the linguistic
and the mathematical level.

The weight vector $\mathbf{w}$ determines which admissibility
dimensions dominate the relaxation dynamics at any given time.
In the sycophancy analysis, high $w_{\mathrm{conv}}$ relative to
$w_{\mathrm{ext}}$ and $w_{\mathrm{adv}}$ produces the pathological
attractor.
In the critical regime, $\mathbf{w}$ is dynamically modulated:
exploration mode increases $w_{\mathrm{sem}}$, verification mode
increases $w_{\mathrm{ext}}$, adversarial mode increases
$w_{\mathrm{adv}}$.

The current paper analyzes only the scalarized regime because the
convergence results of §\ref{sec:convergence}---the Lyapunov
stability argument, the contraction mapping, and the bounded
topology deformation condition---rely on monotonic scalar descent.
The fully vector-valued case requires generalized multi-objective
stability theory, Pareto-admissibility criteria, and potentially
vector Lyapunov functions.
That analysis is deferred to future work.
What the present framework establishes is that the scalar formulation
is a principled special case of the vector structure, not an
approximation of something fundamentally different.
\end{interpretation}

\begin{remark}
In the vector-valued setting, admissibility cannot be defined by a
single threshold $\Phiten(X) < \tau$.
Instead, a configuration is admissible if $\Phiten_k(X) < \tau_k$
for each dimension $k$---a conjunction of component-wise thresholds.
This yields a \emph{rectangular admissibility region} in
$\mathbb{R}^n$, which becomes a convex polytope when the thresholds
are coupled by linear constraints.
The partially ordered structure of admissibility space then reflects
the partial order on $\mathbb{R}^n$ induced by component-wise
dominance: $X \preceq_{\Aspace} Y$ iff $\Phiten_k(X) \leq
\Phiten_k(Y)$ for all $k$.
This is a richer ordering than the scalar case and naturally supports
the analysis of incompatible admissibility pressures---configurations
that improve along one tension dimension while worsening along
another.
\end{remark}

\subsection{Admissibility Failure Metrics}

The admissibility conditions (C1)--(C3) define what it means for a
configuration to be admissible.
But the framework also requires measurable quantities that quantify
\emph{degrees of failure}---not just binary pass/fail judgments.
These metrics make the framework diagnostically useful and provide
candidate empirical targets for future evaluation.

\begin{definition}[Failure Diagnostic Suite]
\label{def:failure}
The following three scalar quantities measure distinct failure modes
of semantic configurations:

\begin{enumerate}
  \item[\textbf{(F1)}] \textbf{Recursive Extension Failure.}
    \[
      E_{\mathrm{ext}}(X)
      \;=\;
      \inf_{\mathcal{E}}\,
      \bigl[\Phiten(\mathcal{E}(X)) - \Phiten(X)\bigr],
    \]
    where the infimum is taken over all admissible extension operators
    $\mathcal{E}$ consistent with $\Cstruct_t$.
    When $E_{\mathrm{ext}}(X) > 0$, every possible continuation
    increases global tension; the configuration is a dead end under
    recursive extension.
    This is the quantitative form of condition (C2) failure.

  \item[\textbf{(F2)}] \textbf{Projection Instability.}
    \[
      E_{\mathrm{proj}}(X^*)
      \;=\;
      \mathrm{Var}_{\ell \,\in\, \mathrm{Lin}(G(X^*))}\,
      \Psi(\ell),
    \]
    the variance of the projection cost $\Psi$ over all linear
    extensions of the stabilized DAG $G(X^*)$.
    High variance indicates that the linearization is highly sensitive
    to small perturbations in the pressure field $\mathcal{P}$---the
    semantic content does not strongly determine a preferred sequential
    form, making the output fragile under paraphrase or translation.

  \item[\textbf{(F3)}] \textbf{Topology Instability.}
    \[
      E_{\Gamma}(t)
      \;=\;
      \left\|\Cstruct_{t+1} - \Cstruct_t\right\|_{\mathrm{F}},
    \]
    the Frobenius norm of the graph adjacency change between successive
    relaxation steps, measuring turbulence in the interaction topology.
    Sustained high $E_{\Gamma}$ signals that the topology is failing to
    converge, which---by the bounded deformation condition
    (Remark~\ref{rem:bounded})---implies that relaxation is not making
    sufficient tension progress to justify the rewiring.
\end{enumerate}

The three metrics combine into an aggregate diagnostic:
\[
  E_{\mathrm{total}}(X, t)
  \;=\;
  \alpha\, E_{\mathrm{ext}}(X)
  \;+\;
  \beta\, E_{\mathrm{proj}}(X^*)
  \;+\;
  \gamma\, E_{\Gamma}(t),
\]
with weighting parameters $\alpha, \beta, \gamma \geq 0$ tunable to
the evaluation regime.
\end{definition}

\begin{interpretation}
The three failure metrics form a coherent diagnostic triad addressing
distinct layers of the SRN architecture.

$E_{\mathrm{ext}}$ is a \emph{forward-looking} measure: it asks
whether the current configuration can be continued.
High $E_{\mathrm{ext}}$ identifies premature stabilization in dead-end
basins---configurations that are locally coherent but globally
exhausted.
This is the formal footprint of the semantic equivalent of a logical
dead end: a line of argument that is locally valid but leads nowhere
further.

$E_{\mathrm{proj}}$ is a \emph{structural} measure: it asks how
robustly the stabilized configuration determines its sequential form.
High $E_{\mathrm{proj}}$ identifies semantic configurations whose
dependency structure is too shallow or too symmetric to strongly
constrain linearization---the SRN analogue of underspecified semantic
content.
This metric would be particularly useful for evaluating ambiguous
paraphrase, stylistic variation, and translation difficulty.

$E_{\Gamma}$ is a \emph{dynamical} measure: it asks whether the
relaxation process is well-behaved.
High sustained $E_{\Gamma}$ identifies oscillatory or turbulent
relaxation dynamics, which correspond computationally to unstable
attention patterns and phenomenologically to circular or incoherent
generation.

The aggregate $E_{\mathrm{total}}$ provides a single scalar
benchmarking target.
A generation system with lower $E_{\mathrm{total}}$ is producing
outputs that are more recursively extensible, more robustly
linearizable, and more dynamically stable than one with higher
$E_{\mathrm{total}}$---properties that correspond directly to
human judgments of text quality, coherence, and fluency.
\end{interpretation}

\subsection{The Recursive Relaxation Operator}

The central dynamical equation of the SRN framework is:
\begin{equation}
  \label{eq:relaxation}
  X_{t+1}
  \;=\;
  \Rop\!\left(X_t,\; -\nabla\Phiten(X_t),\; \Cstruct_t\right),
\end{equation}
where $X_t \in \Aspace$ is the current semantic configuration at relaxation
step $t$, $-\nabla\Phiten(X_t)$ is the descent direction on the tension
functional, $\Cstruct_t$ is the current semantic topology structure, and
$\Rop$ is the \emph{recursive relaxation operator}.

\begin{definition}[Recursive Relaxation Operator]
\label{def:relaxation}
The operator $\Rop : \Aspace \times T\mathcal{X} \times \mathfrak{C} \to \Aspace$
maps the current configuration, the tension gradient, and the
interaction topology to a successor configuration satisfying:
\begin{enumerate}
  \item[(i)] $\Phiten(X_{t+1}) \leq \Phiten(X_t)$ (tension non-increase);
  \item[(ii)] $X_{t+1} \in \Aspace$ (admissibility preservation);
  \item[(iii)] $\Rop$ may invoke rollback: if no forward step satisfies
    (i) and (ii), the operator reverts to a prior stabilization
    checkpoint and explores an alternative relaxation path.
\end{enumerate}
Here $\mathfrak{C}$ denotes the space of semantic topology
configurations.
\end{definition}

\begin{interpretation}
Condition (iii) is the critical departure from standard gradient
descent.
Ordinary gradient flow cannot backtrack; it commits to local descent
and may become trapped in metastable states.
The recursive relaxation operator is explicitly equipped with rollback
capacity: if a locally plausible step leads to a globally inadmissible
configuration, the operator can return to a previous checkpoint.
This is the formal analogue of semantic coherence revision in
cognition---the experience of recognizing that a previously adopted
interpretation was incorrect and revising it in light of later
information.
\end{interpretation}

\begin{remark}
Equation~\eqref{eq:relaxation} is deliberately more general than
standard gradient descent.
The tension functional $\Phiten$ need not be smooth, the constraint
structure $\Cstruct_t$ evolves during relaxation (see
Section~\ref{sec:topology}), and the operator $\Rop$ may perform
non-local updates.
The closest physical analogues are \emph{constraint propagation
algorithms} and \emph{metastable field evolution} in statistical
physics \citep{parisi1988statistical, wilson1974renormalization}.
\end{remark}


\subsection{Operator Stratification}

The recursive relaxation operator $\Rop$ as introduced in
equation~\eqref{eq:relaxation} is deliberately general.
To prevent this generality from collapsing into a universal
placeholder, we explicitly stratify $\Rop$ into three functionally
distinct sub-operators with non-overlapping responsibilities.

\begin{definition}[Stratified Relaxation Operators]
\label{def:stratified}
The full relaxation step decomposes as:
\[
  \Rop(X_t, -\nabla\Phiten(X_t), \Cstruct_t)
  \;=\;
  \Rop^{(3)} \circ \Rop^{(2)} \circ \Rop^{(1)},
\]
where:
\begin{enumerate}
  \item[(i)] $\Rop^{(1)}$ is the \textbf{local relaxation operator}:
    performs small admissibility-preserving gradient steps on
    individual regions $X_i \in \mathcal{F}_t$ (the active instability
    frontier; see Section~\ref{subsec:frontier}).
    $\Rop^{(1)}$ does not modify topology or invoke rollback.

  \item[(ii)] $\Rop^{(2)}$ is the \textbf{topological restructuring
    operator}: applies the graph rewiring step $\Gop(\Cstruct_t, X_t,
    \nabla\Phiten(X_t))$, updating coupling weights and adjacency
    structure subject to the bounded deformation constraint.
    $\Rop^{(2)}$ does not modify region content or invoke rollback.

  \item[(iii)] $\Rop^{(3)}$ is the \textbf{verification and rollback
    operator}: applies the hierarchical verification field
    $\mathcal{V}$ to the output of $\Rop^{(2)} \circ \Rop^{(1)}$.
    If verification passes at all active scales, the step is accepted.
    If verification fails at any scale, $\Rop^{(3)}$ reverts the
    configuration to the most recent valid checkpoint and marks the
    failed path as explored.
\end{enumerate}
\end{definition}

\begin{remark}
This stratification resolves a potential ambiguity in the original
formulation: the single operator $\Rop$ appeared capable of performing
arbitrarily many distinct operations simultaneously.
The stratified form makes clear that local relaxation, topology
adaptation, and verification are sequential sub-operations within each
step, operating over non-overlapping aspects of the system state.
This separation also clarifies the computational budget: $\Rop^{(1)}$
is the cheapest sub-operator (sparse local updates); $\Rop^{(2)}$ is
moderate (graph rewiring); $\Rop^{(3)}$ is the most expensive in
the worst case (full hierarchical verification) but is expected to
terminate quickly for well-stabilized regions.
\end{remark}


\subsection{The Active Instability Frontier and Locality Recovery}
\label{subsec:frontier}

A natural objection to the SRN framework is computational: if
generation requires minimizing a global tension functional over a
high-dimensional semantic configuration, the architecture appears to
require globally expensive computation at every step---far more costly
than the local attention operations of a transformer.

This objection misreads the dynamics.
Locality is not abandoned in SRNs; it is \emph{recovered dynamically}
through instability concentration.

At any relaxation step $t$, the overwhelming majority of semantic
regions will have already achieved local stability: their instability
measures satisfy $S_i(X_t) < \tau$.
These regions require no further relaxation computation.
High-cost updates concentrate on the \emph{active instability
frontier}:

\begin{definition}[Active Instability Frontier]
\label{def:frontier}
The \emph{active instability frontier} at step $t$ is the set of
regions that have not yet achieved local stability:
\[
  \mathcal{F}_t
  \;=\;
  \left\{\, X_i \;:\; S_i(X_t) > \tau_{\mathrm{active}} \,\right\},
\]
for an active threshold $\tau_{\mathrm{active}} \leq \tau$.
The local relaxation operator $\Rop^{(1)}$ applies exclusively to
regions in $\mathcal{F}_t$ and their immediate topological neighbors
in $\Cstruct_t$.
\end{definition}

\begin{principle}[Locality Recovery]
\label{prin:locality}
Under the active instability frontier, the per-step computational cost
of $\Rop^{(1)}$ scales with $|\mathcal{F}_t|$ rather than with the
total number of semantic regions $|V|$.
As relaxation proceeds and regions stabilize, $|\mathcal{F}_t|$
decreases monotonically (subject to rollback events, which may
temporarily re-activate previously stabilized regions).
In the limit, $|\mathcal{F}_t| \to 0$ as $X_t \to X^*$.
\end{principle}

\begin{interpretation}
This locality recovery has two important consequences.

First, it makes the computational cost of SRN generation
\emph{front-loaded}: expensive diffuse computation early in relaxation
(when instability is widespread) gives way to cheap sparse computation
as the configuration approaches its admissibility basin.
This is the opposite of autoregressive generation, whose per-token
cost is roughly constant throughout the sequence.
For tasks where early-stage semantic commitment is cheap and late-stage
refinement is expensive---long-form generation, complex argumentation,
multi-document synthesis---the SRN profile may be computationally
advantageous.

Second, this account has a direct cognitive parallel.
Human cognition does not continuously recompute semantic structure
across the full breadth of active representations.
Attention---in both the cognitive and colloquial sense---concentrates
around regions of unresolved tension: the word that doesn't fit, the
argument that doesn't follow, the reference that hasn't been resolved.
The active instability frontier is a formal analogue of this
attentional focusing: expensive processing concentrates precisely where
instability is highest, withdrawing from regions that have reached
local coherence.

In practical terms, a sparse SRN implementation would maintain
$\mathcal{F}_t$ as an explicit priority queue ordered by $S_i$,
processing only regions above threshold and propagating tension updates
to their topological neighbors.
This yields an architecture whose computational behavior is much closer
to sparse attention mechanisms than to full-field global optimization.
\end{interpretation}

\subsection{Emergent Tokens as Crystallization}

A crucial consequence of the SRN framework is a reconceptualization of
what tokens are.

In autoregressive models, tokens are primitive input and output units.
They are the atoms of generation: everything else is built from them.

In SRNs, tokens are \emph{emergent crystallizations}.

When a sub-region $X_i$ achieves a local instability measure below the
coherence threshold---$S_i(X_t) < \tau$---and no further relaxation
step produces a meaningful reduction in $S_i$, the region has
\emph{stabilized}.
A stabilized region can be projected into the communication channel as
a token or token sequence.
This projection is not generation in the autoregressive sense; it is
\emph{readout}: the conversion of a stabilized admissibility region
into a linear symbolic trace.

The model does not ask: \emph{which token comes next?}

It asks: \emph{which regions of semantic configuration space have
stabilized sufficiently to be communicated?}

This distinction is not merely semantic (in the colloquial sense).
It has concrete architectural consequences: multiple distant regions
may stabilize simultaneously; stabilization may be partial and require
further relaxation; and the order in which regions collapse into tokens
need not match the left-to-right order of the eventual communicative
trace.

% ==============================================================
\section{Dynamic Semantic Topology Formation}
\label{sec:topology}

\subsection{The Fixed-Topology Problem}

The transformer attention mechanism computes interactions between
all pairs of positions with weights determined by the softmax of
scaled dot products:
\[
  \mathrm{Attn}(Q, K, V)
  \;=\;
  \mathrm{softmax}\!\left(\frac{QK^\top}{\sqrt{d_k}}\right) V.
\]
This is an extraordinarily powerful mechanism.
But its interaction topology is determined by sequence position and
content---it is not determined by \emph{semantic constraint structure}.

Two positions that are semantically tightly coupled but positionally
distant will compete for attention weight against all other positions.
Two positions that are positionally adjacent but semantically
uncoupled will exchange attention weight regardless.

The topology of interaction is, in a precise sense,
\emph{sequence-indexed}: it is downstream of position, not upstream of
meaning.

\subsection{The Topology Evolution Equation}

The SRN framework proposes that the semantic interaction topology
should itself be a dynamical object, evolving in response to the
developing constraint structure.

The \emph{topology evolution equation} is:
\begin{equation}
  \label{eq:topology}
  \Cstruct_{t+1}
  \;=\;
  \Gop\!\left(\Cstruct_t,\; X_t,\; \nabla\Phiten(X_t)\right),
\end{equation}
where $\Cstruct_t$ is the current semantic adjacency and coupling structure
(an evolving graph over semantic regions), $\Gop$ is the
\emph{topology reconfiguration operator}, and the graph evolves
according to constraint gradients and stabilization pressures.

\begin{definition}[Topology Reconfiguration Operator]
\label{def:topology}
The operator $\Gop : \mathfrak{C} \times \Aspace \times T\mathcal{X} \to
\mathfrak{C}$ updates the semantic interaction graph according to the
following principles:
\begin{enumerate}
  \item[(i)] \textbf{Constraint coupling.}
    An edge $(i,j)$ in $\Cstruct_{t+1}$ has weight $\omega_{ij}^{t+1}$
    proportional to the semantic tension $T(i,j)$ at step $t$:
    high-tension region pairs become more strongly coupled.
  \item[(ii)] \textbf{Stabilization decoupling.}
    Once a region $X_i$ achieves local stability ($S_i < \tau$),
    its coupling to unstabilized regions is reduced, allowing
    the relaxation dynamics to focus computational resources on
    remaining instability.
  \item[(iii)] \textbf{Proximity-independent adjacency.}
    Two regions $X_i$ and $X_j$ that are positionally distant may
    become \emph{semantically adjacent} in $\Cstruct_{t+1}$ if their
    mutual constraint tension $T(i,j)$ exceeds a coupling threshold.
\end{enumerate}
\end{definition}

\begin{interpretation}
Principle (iii) is the most architecturally radical element of the
framework.
In a transformer, the attention graph is approximately sequence-local
(modified by content, but bounded by the sequence structure).
In an SRN, the semantic interaction graph can be non-local: two
distant textual regions become temporarily adjacent if they exhibit
strong mutual constraint tension.
The graph is not a fixed scaffold imposed on the generation process; it
is an emergent structure that self-organizes according to the
constraint dynamics of the unfolding configuration.
This is what we mean when we say that SRNs operate over a dynamically
evolving admissibility topology rather than a static sequence manifold.
\end{interpretation}


\begin{remark}[Conservation and Bounded Topology Deformation]
\label{rem:bounded}
The topology reconfiguration operator $\Gop$ as defined above is
relatively unconstrained: given admissibility preservation, the graph
may reorganize freely at each step.
This freedom is conceptually important---it is what allows SRNs to
form non-local semantic adjacencies and to escape the sequence-indexed
topology of transformers.

However, unlimited topology plasticity creates a potential instability.
If the graph reorganizes faster than the relaxation dynamics can
dissipate tension, the system may fail to converge: new edges
introduce new tension sources before old ones are resolved.

We therefore impose a \emph{bounded topology deformation} condition:
\[
  \left\|\Cstruct_{t+1} - \Cstruct_t\right\|_{\mathrm{F}}
  \;\leq\;
  \eta_{\Gamma} \cdot \left|\Phiten(X_t) - \Phiten(X_{t-1})\right|,
\]
where $\|\cdot\|_{\mathrm{F}}$ is the Frobenius norm on the graph
adjacency structure and $\eta_{\Gamma} > 0$ is a topology plasticity
rate.
This condition ties the rate of graph reorganization to the rate of
tension reduction: the topology may only restructure as quickly as
the relaxation is making progress.

Future implementations may strengthen this to a conservation-like
condition, requiring that total coupling weight $\sum_{i,j}
\omega_{ij}^t$ is approximately conserved under $\Gop$, so that
topology evolution redistributes rather than amplifies coupling.
This would prevent the semantic interaction graph from becoming
arbitrarily dense or sparse under repeated reconfiguration, and would
provide a natural entropy budget on topology evolution analogous to
the entropy constraints in the RSVP framework.
\end{remark}

\begin{proposition}
\label{prop:topology}
Under the joint dynamics of equations~\eqref{eq:relaxation}
and~\eqref{eq:topology}, the semantic topology $\Cstruct_t$ converges
toward a structure in which strongly coupled regions correspond
precisely to regions of high mutual semantic dependency in the
stabilized configuration $X^*$.
That is, the emergent topology encodes the semantic structure of the
output rather than its sequential order.
\end{proposition}

\subsection{Adaptive Semantic Neighborhoods}

A consequence of dynamic topology formation is that semantic
neighborhoods $\Nhd_i$ are not fixed.
At each relaxation step, the neighborhood of region $X_i$ consists of
the set of regions currently coupled to it in $\Cstruct_t$:
\[
  \Nhd_i^t \;=\; \left\{ j \neq i \;:\; \omega_{ij}^t > 0 \right\}.
\]
As the topology evolves, neighborhoods contract and expand.
A region that was previously in isolation (low coupling with all
others) may later be drawn into close constraint relationship with a
distant region when both fail to stabilize independently.

This adaptive neighborhood structure replaces both the fixed attention
window of standard transformers and the entropy-based patching of
BLT-style architectures.
The organizational principle is not position, not entropy, but
\emph{constraint coupling}.

% ==============================================================
\section{Stabilization, Verification, and Hallucination}
\label{sec:verification}

\subsection{Hierarchical Verification Operators}

BLT introduces a dynamic verification mechanism (BLT-DV) that checks
local patch predictions against a latent global model.
The SRN framework generalizes this intuition significantly.

\begin{definition}[Verification Field]
\label{def:verification}
A \emph{hierarchical verification field} is a family of operators
$\mathcal{V} = \{ V_k \}_{k=1}^K$ where each $V_k$ checks whether
local relaxation at scale $k$ preserves global admissibility
invariants at scale $k+1$:
\[
  V_k(X_{t+1}, \Cstruct_{t+1}) \;=\;
  \begin{cases}
    \checkmark & \text{if } X_{t+1} \in \Aspace \text{ at scale } k+1, \\
    \text{rollback to } X_{t'} & \text{otherwise}.
  \end{cases}
\]
\end{definition}

The verification field enables:
\begin{itemize}
  \item \textbf{Local speculative collapse}: a sub-region may
    tentatively crystallize into tokens while remaining subject to
    global verification.
  \item \textbf{Global semantic rollback}: if a local crystallization
    violates a higher-scale admissibility invariant, the operator
    reinstates a prior checkpoint.
  \item \textbf{Hierarchical reconciliation}: tension at scale $k$
    is resolved by adjusting constraint coupling at scale $k-1$.
  \item \textbf{Asynchronous stabilization}: different regions may
    stabilize at different rates, proceeding through verification
    independently.
\end{itemize}

\subsection{Hallucination as Topological Instability}

We now return to the phenomenon of hallucination with the conceptual
apparatus developed above.

\begin{principle}[Hallucination as Premature Collapse]
\label{prin:hallucination}
Hallucination in an SRN occurs when a local region $X_i$ achieves
apparent stability---$S_i(X_t) < \tau$---and collapses into tokens
before the hierarchical verification field $\mathcal{V}$ has confirmed
global admissibility.
The region crystallizes prematurely, emitting tokens that are locally
coherent but globally inadmissible.
\end{principle}

\begin{interpretation}
This reinterpretation of hallucination is explanatorily deeper than the
standard account.
Under sequential ontology, hallucination is: the model assigned high
probability to an incorrect token.
Under SRN, hallucination is: the stabilization dynamics allowed a
local collapse without verifying recursive extensibility at higher
scales.
This immediately suggests a structural remedy: the verification
field $\mathcal{V}$ is precisely the mechanism that prevents premature
collapse by requiring that local stabilization survive hierarchical
consistency checks before crystallizing.

The analogy with physical systems is instructive.
Crystal growth exhibits analogous phenomena: a local lattice
arrangement may appear energetically stable transiently but later
generate structural incompatibilities that propagate as defects.
Hallucination is the cognitive analogue of crystallographic defect
propagation---a local structure that violated global lattice
admissibility before being locked in.
\end{interpretation}

% ==============================================================
\section{Emergent Sequentiality and Linguistic Projection}
\label{sec:emergent}

\subsection{The Inversion}

We are now in a position to state the deepest conceptual claim of the
paper.

Sequential language is not the substrate of cognition.
It is the \emph{observable residue} of recursive constraint
stabilization processes operating over high-dimensional admissibility
fields.

This is an inversion of the standard picture.

The standard picture says: cognition processes sequences; language is a
sequence; language models are cognitive models.

The SRN picture says: cognition is distributed recursive stabilization;
language generation is the projection of stabilized admissibility
regions onto the one-dimensional communication channel; sequential text
is what stabilization looks like from outside.

\begin{principle}[Emergent Sequentiality]
\label{prin:sequentiality}
Let $X^* \in \Aspace$ be a fully stabilized semantic configuration.
The sequential output $x = (x_1, x_2, \ldots, x_n)$ is a
\emph{readout projection}:
\[
  x \;=\; \pi\!\left(X^*\right),
\]
where $\pi : \Aspace \to \Sigma^*$ is a projection from admissibility space
onto the space of symbol sequences $\Sigma^*$.
The projection $\pi$ selects a linear ordering of the stabilized
regions of $X^*$ consistent with communicative conventions
(grammar, discourse structure, pragmatic context).
The ordering is \emph{induced} by communicative constraints, not by
the stabilization process itself.
\end{principle}

\begin{interpretation}
This principle has several consequences.

First, \emph{multiple linear orderings may project from the same
stabilized configuration}.
The same underlying semantic structure may be expressible in different
sequential forms (different word orders, different discourse
organizations) without any change to the admissibility structure
itself.

Second, \emph{the sequential order is not the computationally primary
object}.
It is a post-hoc linearization of a non-sequential stabilization
process.
Language models that treat sequential order as computationally primary
are, in effect, learning to reconstruct the projection while remaining
blind to the underlying admissibility geometry.

Third, \emph{the difficulty of many NLP tasks---coreference,
long-range agreement, discourse coherence---has a unified explanation}:
these tasks require recovering non-local admissibility relationships
that were collapsed by the linearization projection.
Sequential models must reconstruct admissibility structure from its
projection, which is informationally incomplete.

Fourth, this account connects naturally to human language evolution.
The communication channel constraint---speech and writing are
one-dimensional---would have exerted pressure for a stable projection
convention.
Grammar and syntax, on this view, are \emph{projection conventions}:
stable mappings from high-dimensional semantic admissibility structure
onto one-dimensional communicative sequences.
\end{interpretation}


\subsection{The Minimal SRN}
\label{subsec:minimal}

Before describing the full SRN architecture, it is important to
distinguish the \emph{irreducible core} of the framework from the
broader ecosystem of operators and mechanisms developed in this paper.
RSVP, TARTAN, and CLIO provide theoretical depth and implementation
infrastructure, but they are not required for the essential SRN
dynamics.

\begin{principle}[Minimal SRN]
\label{prin:minimal}
A minimal SRN requires exactly five components:
\begin{enumerate}
  \item[\textbf{(M1)}] A \textbf{semantic configuration space}
    $\mathcal{X}$ equipped with a partial order and a notion of local
    constraint coherence.
  \item[\textbf{(M2)}] A \textbf{tension functional} $\Phiten :
    \mathcal{X} \to \mathbb{R}_{\geq 0}$ measuring total semantic
    incompatibility.
  \item[\textbf{(M3)}] A \textbf{relaxation operator} $\Rop$ that
    reduces $\Phiten$ at each step while preserving admissibility,
    with rollback capacity for failed steps.
  \item[\textbf{(M4)}] A \textbf{topology adaptation mechanism}
    $\Gop$ that reorganizes the interaction graph in response to
    constraint gradients.
  \item[\textbf{(M5)}] A \textbf{projection operator} $\pi$ that
    converts a stabilized configuration into a sequential output by
    selecting a linear extension of the residual dependency structure.
\end{enumerate}
Everything else in this paper---the hierarchical verification field,
the communicative pressure field decomposition, the failure diagnostic
suite, the TARTAN memory reservoir---constitutes scalable refinement
of this minimal core.
\end{principle}

\begin{interpretation}
This invariant specification has a practical consequence: a minimal
SRN can be implemented at small scale without any of the additional
machinery.
A semantic graph with explicit coupling weights, a tension-minimizing
relaxation rule, a graph rewiring step, a simple rollback mechanism,
and a topological sort for output constitute a functional, testable
SRN.

This is the architecture of a constrained graph relaxation system with
rollback checkpoints---substantially simpler than a transformer, and
directly testable on tasks where sequential commitment is structurally
costly: garden-path disambiguation, long-distance agreement,
coreference resolution, and discourse repair.
Such a system would not generate novel text.
But it would operationalize the paper's central empirical claim: that
dynamically stabilized semantic graphs with rollback resolve
ambiguity more robustly than left-to-right sequential commitment.
That demonstration alone would constitute a meaningful empirical
contribution.
\end{interpretation}

\subsection{SRN Architecturally}

Translating the above into computational architecture, an SRN consists
of the following principal components:

\begin{enumerate}
  \item \textbf{Semantic Constraint Field Encoder}: Maps input
    (text, other modalities, or prior context) into an initial
    semantic configuration $X_0 \in \mathcal{X}$, not necessarily
    in $\Aspace$ initially.

  \item \textbf{Admissibility Gradient Network}: Computes the
    tension gradient $\nabla\Phiten(X_t)$ and local instability measures
    $S_i$ with respect to the current topology $\Cstruct_t$.

  \item \textbf{Recursive Relaxation Operator}: Implements
    equation~\eqref{eq:relaxation}, including rollback capacity and
    admissibility preservation.

  \item \textbf{Topology Reconfiguration Layer}: Implements
    equation~\eqref{eq:topology}, maintaining and evolving the
    semantic interaction graph $\Cstruct_t$.

  \item \textbf{Hierarchical Verification Field}: Implements the
    multi-scale verification operators $\mathcal{V}$, preventing
    premature crystallization.

  \item \textbf{Persistence Memory Reservoir}: Stores stabilized
    admissibility regions and unresolved tension structures rather than
    previous tokens.
    Memory is \emph{geometric}: it encodes constraint configurations,
    not symbol sequences.

  \item \textbf{Readout Projection}: Implements $\pi$, converting
    stabilized regions into sequential token output when appropriate
    stability thresholds are met.
\end{enumerate}

This architecture replaces the key-value cache of transformer-based
systems not with a larger cache but with a geometrically organized
persistence reservoir.
The reservoir stores not \emph{what was said} but \emph{what has been
stabilized}---the constraint relationships that govern future
admissible extension.


\subsection{Worked Example: The Garden-Path Sentence}

The theoretical machinery developed above becomes concrete when applied
to a canonical failure case for sequential generation: the garden-path
sentence.

Consider the sentence:

\begin{center}
  \textit{The horse raced past the barn fell.}
\end{center}

This sentence is grammatically well-formed but produces systematic
processing difficulty in human readers and catastrophic failure in
autoregressive language models.
The difficulty is structural: the initial subsequence \textit{the
horse raced past the barn} supports a strong main-clause reading
(``the horse raced''), but the final word \textit{fell} requires
reanalysis of \textit{raced} as a reduced relative clause modifier
(``the horse [that was] raced past the barn'').

\subsubsection*{Autoregressive Failure}

Under an autoregressive model, generation or comprehension of this
sentence proceeds token by token:

\medskip
\noindent
$x_1 = $ \textit{The}, $x_2 = $ \textit{horse}, $x_3 = $
\textit{raced}, \ldots

\medskip
\noindent
At position $x_3$, the model assigns high conditional probability to
\textit{raced} as the main verb.
This commitment is irrevocable: the hidden state at step $t=3$
encodes \textit{raced} as a tensed main clause predicate.
When \textit{fell} arrives at position $x_8$, the model has no
mechanism to revise the role of \textit{raced}.
The result is either a sharp drop in likelihood (if the model is
scoring), an inability to parse the sentence, or an erroneous
completion that avoids \textit{fell} entirely.

This is not a failure of attention or memory.
The transformer can, in principle, attend to all previous positions
when processing \textit{fell}.
The failure is ontological: the sequential commitment to a particular
semantic role assignment for \textit{raced} cannot be undone.

\subsubsection*{SRN Relaxation Trace}

Under an SRN, generation proceeds through stabilization of semantic
regions rather than commitment to sequential tokens.
The following trace is schematic, indicating the qualitative evolution
of the key regions and tensions.

\medskip
\noindent
\textbf{Initialization} ($t = 0$):
The configuration $X_0$ encodes the full semantic content---agent,
event, modification, main predicate---as an unstabilized field.
Local instability is high throughout: $S_i(X_0) \gg \tau$ for all
regions. The topology $\Cstruct_0$ is dense, with all regions weakly
coupled.

\medskip
\noindent
\textbf{Early relaxation} ($t = 1, 2$):
The agent region $X_{\mathrm{agent}}$ (encoding
\textsc{horse}) stabilizes rapidly: $S_{\mathrm{agent}} < \tau$.
Two candidate event regions form: $X_{\mathrm{main}}$ (encoding
\textsc{horse} as agent of a main-clause event) and $X_{\mathrm{rel}}$
(encoding \textsc{horse} as argument of a relative-clause event).
Both are transiently admissible.
The tension $T(\mathrm{main}, \mathrm{rel})$ is high: the two
configurations are mutually incompatible under (C1).

\medskip
\noindent
\textbf{Premature collapse} ($t = 3$):
If the verification field $\mathcal{V}$ does not intervene, the
relaxation dynamics may favour $X_{\mathrm{main}}$ (higher local
probability, lower local instability at this step).
$X_{\mathrm{main}}$ crystallizes tentatively.
This corresponds to the garden-path reading: the model has committed
to the main-clause interpretation.

\medskip
\noindent
\textbf{Constraint arrival and tension spike} ($t = 4$):
The main-predicate region $X_{\mathrm{fell}}$ (encoding
\textsc{fell} as a tensed verb) becomes active.
Under the main-clause reading, $X_{\mathrm{fell}}$ generates
catastrophic tension with $X_{\mathrm{main}}$: two tensed predicates
compete for the same clausal slot.
The global tension functional spikes:
$\Phiten(X_t) \gg \Phiten(X_{t-1})$.

\medskip
\noindent
\textbf{Rollback} ($t = 5$):
The recursive relaxation operator $\Rop$, governed by condition (R3),
detects the tension spike and activates rollback to the most recent
checkpoint preceding crystallization of $X_{\mathrm{main}}$---that
is, the state at $t = 2$ before the premature collapse.
The tentative crystallization of $X_{\mathrm{main}}$ is revoked.

\medskip
\noindent
\textbf{Topology reorganization} ($t = 5, 6$):
Following rollback, the topology reconfiguration operator $\Gop$
increases the coupling between $X_{\mathrm{rel}}$ and
$X_{\mathrm{fell}}$: these two regions exhibit strong mutual
constraint compatibility (relative-clause event $+$ main predicate is
a coherent configuration).
The edge $(\mathrm{rel}, \mathrm{fell})$ is strengthened in
$\Cstruct_5$.
Simultaneously, $T(\mathrm{rel}, \mathrm{fell})$ decreases toward
zero as the two regions co-stabilize.

\medskip
\noindent
\textbf{Convergence} ($t = 7$):
$X_{\mathrm{rel}}$ and $X_{\mathrm{fell}}$ stabilize jointly.
The configuration $X^*$ encodes the correct reduced-relative reading.
$\Phiten(X^*) \approx 0$.

\medskip
\noindent
\textbf{Projection} ($t = 8$):
The semantic constraint DAG $G(X^*)$ has $X_{\mathrm{agent}}$ and
$X_{\mathrm{rel}}$ preceding $X_{\mathrm{fell}}$ in the dependency
order.
The communicative pressure field $\mathcal{P}$ selects the linear
extension:
\[
  \pi(X^*)
  \;=\;
  \langle\,
  \textit{The horse raced past the barn fell.}\,
  \rangle,
\]
correctly linearizing the reduced-relative structure.

\begin{interpretation}
Several features of this trace are worth highlighting.

The rollback at $t = 5$ is the key operation that autoregressive
models cannot perform.
It requires that local crystallizations remain \emph{provisional}
until hierarchical verification is satisfied---exactly the function of
the verification field $\mathcal{V}$.

The topology reorganization at $t = 5$--$6$ illustrates proximity-independent
adjacency: $X_{\mathrm{rel}}$ and $X_{\mathrm{fell}}$ are not
adjacent in any positional sense (they correspond to non-adjacent
spans), but their mutual constraint compatibility draws them into close
coupling under $\Gop$.

The projection at $t = 8$ is distinct from the stabilization process.
The SRN does not generate \textit{The horse raced past the barn fell}
sequentially.
It stabilizes the semantic configuration and then linearizes it.
The sequential output is the residue of a non-sequential process.

Finally, note that the garden-path effect in human processing
corresponds to the tension spike and rollback at steps $4$--$5$.
Human readers experience processing difficulty at \textit{fell}
precisely because the cognitive system must revoke a committed
interpretation---a rollback operation---and restructure its dependency
representation.
SRNs model this naturally; autoregressive models have no corresponding
mechanism.
\end{interpretation}


% ==============================================================
\section{SRNs as Generalized Energy-Based Systems}
\label{sec:ebm}

The SRN framework will immediately remind many readers of
energy-based models (EBMs), a well-established family of
architectures in which generation corresponds to minimizing a learned
energy function \citep{lecun2006ebm, hopfield1982neural,
ackley1985boltzmann}.
This resemblance is genuine and should be acknowledged directly.
SRNs \emph{are} a species of energy-based system.
But they introduce four structural departures that collectively define
a distinct computational regime.

\subsection{The EBM Baseline}

A standard EBM defines an energy function $E_\theta(x)$ over the
output space and generates outputs by finding low-energy
configurations:
\[
  x^* \;=\; \arg\min_{x} E_\theta(x).
\]
The energy is typically a fixed function of the output and a condition,
with a static interaction structure.
Hopfield networks \citep{hopfield1982neural} and Boltzmann machines
\citep{ackley1985boltzmann} are discrete EBMs; continuous EBMs
\citep{lecun2006ebm} extend this to structured output spaces.
Harmony Theory \citep{smolensky1986harmony} applies EBM-style
optimization specifically to linguistic representations.

SRNs agree with EBMs that generation is optimization of an objective
functional rather than sequential prediction.
The global tension functional $\Phiten(X)$ is the SRN analogue of the
energy $E_\theta(x)$.

\subsection{Four Departures from Standard EBMs}

\begin{enumerate}
  \item \textbf{Dynamic interaction topology.}
    In standard EBMs, the interaction structure between variables is
    fixed by the model architecture.
    The SRN interaction topology $\Cstruct_t$ evolves during
    optimization under the operator $\Gop$, allowing the model to
    reorganize which variables are coupled based on emerging constraint
    structure.
    This is not a feature of any standard EBM formulation.

  \item \textbf{Recursive extensibility as a constraint.}
    Standard EBMs optimize energy over the full output at once or
    iteratively refine a fixed-dimensional configuration.
    SRNs impose the additional constraint that admissible configurations
    must be recursively extensible: the system evaluates not just
    whether the current configuration is low-energy, but whether future
    extensions can maintain low energy.
    Condition (C2) of Definition~\ref{def:admissibility} has no direct
    analogue in standard EBM frameworks.

  \item \textbf{Hierarchical rollback verification.}
    EBMs typically commit to descending the energy landscape monotonically
    (with the exception of simulated annealing variants, which introduce
    stochastic acceptance of uphill steps).
    The SRN verification hierarchy $\mathcal{V}$ introduces structured
    rollback: violations of admissibility at higher scales trigger
    reversion to explicit checkpoints, not stochastic re-sampling.
    This provides a deterministic mechanism for escaping metastable
    traps that is absent from standard EBMs.

  \item \textbf{Projection-field linearization.}
    Standard EBMs produce outputs as configurations over a fixed output
    space.
    SRNs produce outputs through a two-stage process: stabilization of
    a non-sequential configuration, followed by projection via the
    communicative pressure field $\mathcal{P}$.
    The sequential output is not the object being optimized; it is the
    residue of an optimization performed over a higher-dimensional
    admissibility structure.
    This projection stage has no analogue in any standard EBM.
\end{enumerate}

\subsection{Relationship to Diffusion and Predictive Coding}

Diffusion-based language models \citep{ho2020ddpm, austin2021discrete}
share with SRNs the intuition that generation is an iterative
refinement process rather than sequential emission.
The denoising process in diffusion models can be read as a form of
tension reduction: corrupted configurations are driven toward coherent
ones by a learned score function.
However, diffusion operates over a fixed sequence-indexed state space
and recovers an ordered output; it does not reorganize its interaction
topology or impose recursive extensibility constraints.

Predictive coding networks \citep{friston2010free, friston2006free}
are perhaps the closest existing architecture to SRNs in their
computational philosophy: hierarchical error minimization through
iterative message-passing across levels.
SRNs can be understood as a generalization of predictive coding in
which: the prediction error is replaced by semantic tension, the
fixed hierarchical structure is replaced by a dynamically evolving
topology, and the output layer is replaced by a projection operator
over a stabilized dependency DAG.

Graph neural networks with iterative message-passing
\citep{goodfellow2016deep} provide the most natural implementation
substrate for the local relaxation operator $\Rop^{(1)}$, since the
core computation is neighborhood-aggregated message passing over a
weighted graph---exactly the structure of the SRN interaction topology
$\Cstruct_t$.


A deeper theoretical precedent for the SRN claim that sequential and
parallel computational structures are not primitively distinct comes
from Fant's \emph{Computer Science Reconsidered}
\citep{fant2007computer}.
Fant argues that the apparent complexity of concurrency is not
intrinsic but is an artifact of a conceptual framework inherited from
mathematics that privileges sequential, algorithm-based expression.
His central contribution is the \emph{invocation model} and its
associated NULL Convention Logic: a three-valued signal system in
which every wire carries not just \textsc{true} or \textsc{false}
but also a third state, \textsc{null}, signifying ``not yet data.''
In this framework, completion of computation is expressed
intrinsically---each function asserts \textsc{null} between data
wavefronts, so that the downstream network can detect unambiguously
when a new valid result has arrived and when the system is between
computations---without requiring a global clock to impose sequential
discipline.
The critical result is that any concurrent network of functions
satisfying this completeness criterion can be partitioned and mapped
onto sequential, parallel, or hybrid implementations without changing
the expression of the underlying computation: the network is
expressed once and mapped to any available coordination protocol.
Sequentiality, on Fant's account, is therefore not a primitive
property of computation but an implementation-level coordination
convention imposed on a fundamentally concurrent and distributed
underlying process.

This is a direct computational antecedent to the SRN claim that
sequential language is not a primitive substrate but the projection
residue of distributed stabilization.
Both frameworks locate sequentiality as downstream of a richer,
non-sequential underlying process: Fant's data wavefronts propagating
through concurrent networks of threshold operators correspond
structurally to SRN stabilization fronts propagating through the
active instability frontier $\mathcal{F}_t$.
The NULL state---the signal of incompletion between wavefronts---
corresponds to the unstabilized admissibility state of semantic
regions that have not yet crystallized.
And the completeness criterion of NULL Convention Logic, which
determines when a function's output is valid, corresponds to the
local coherence condition (C1) of Definition~\ref{def:admissibility},
which determines when a semantic region has stabilized sufficiently
to participate in projection.
The coordination mechanism that converts distributed concurrent
completion into apparent sequential order is, in Fant, the
alternating data/null wavefront protocol; in SRNs, it is the
communicative pressure field $\mathcal{P}$ that selects a linear
extension of the stabilized dependency DAG $G(X^*)$.

The SRN framework is therefore best understood not as a wholly novel
architecture but as a \emph{geometric reinterpretation and
unification} of these converging research trajectories: the
optimization orientation of EBMs, the iterative refinement of
diffusion, the hierarchical error-correction of predictive coding, and
the graph-structured computation of GNNs---organized around a common
principle of constraint stabilization over dynamically evolving
admissibility topology.

% ==============================================================
\section{Relationship to RSVP, TARTAN, and CLIO}
\label{sec:rsvp}

\subsection{RSVP as the Foundational Ontology}

The Relativistic Scalar-Vector Plenum (RSVP) framework provides the
broadest ontological context for the present work.
RSVP models physical and semantic reality as a field system in which
scalar (entropic), vector (directional), and plenum (substrate)
components evolve under coupled dynamics.
Admissibility, in RSVP terms, corresponds to configurations that
maintain accessibility under entropy-bounded persistence: a region of
state space is accessible if it can be reached from the current state
via transformations that do not exceed local entropy constraints.

The SRN framework instantiates this general ontology in the specific
domain of linguistic and cognitive generation.
Admissibility space in SRNs is the semantic analogue of the RSVP
accessibility manifold: a configuration is semantically admissible if
it can be recursively extended without catastrophic entropy increase
in the semantic tension functional.

Crucially, the SRN framework is designed to be comprehensible
independently of RSVP.
The admissibility conditions (C1)-(C3) of
Definition~\ref{def:admissibility} are stated in terms that do not
require commitment to RSVP field ontology.
RSVP provides interpretive depth; it is not a prerequisite.

\subsection{TARTAN as Multiscale Memory}

The Trajectory-Aware Recursive Tiling with Annotated Noise (TARTAN)
framework provides the multiscale memory and trajectory buffering
infrastructure that an SRN implementation would require.

TARTAN's hierarchical tiling structure maps naturally onto the SRN's
multi-scale stabilization: different tiling scales correspond to
different levels of the verification hierarchy $\mathcal{V}$.
Stabilized regions at fine scales contribute to coarser-scale
admissibility assessments.
The trajectory buffering mechanism provides the checkpoint storage
necessary for rollback operations in the recursive relaxation operator.

\subsection{CLIO as Recursive Optimization}

The Constraint-Limited Inference and Optimization (CLIO) framework
provides the recursive optimization and verification dynamics that
implement the SRN's stabilization processes.

CLIO's recursive verification structure corresponds directly to the
hierarchical verification field $\mathcal{V}$ of
Definition~\ref{def:verification}.
CLIO's constraint-limited inference mechanism provides the operational
semantics for the admissibility gradient network.

Together, RSVP, TARTAN, and CLIO constitute the broader framework
ecosystem within which SRNs are the linguistic-cognitive specialization:
\begin{itemize}
  \item RSVP: general admissibility-field ontology,
  \item TARTAN: multiscale recursive memory and trajectory management,
  \item CLIO: recursive optimization and stabilization verification,
  \item SRN: linguistic generation as constraint stabilization.
\end{itemize}

% ==============================================================
\section{Implications for AI, Cognition, and Physics}
\label{sec:implications}

\subsection{For Artificial Intelligence}

The SRN framework suggests a reorientation of the language model
research agenda.

Evaluation should move beyond perplexity and token-level accuracy
toward metrics that capture \emph{admissibility}: semantic coherence
under perturbation, recursive extensibility, and resistance to
premature local collapse.
These are not the same as factual accuracy, though the two are related.

Architecture design should explore mechanisms for dynamic topology
formation, persistence reservoirs over constraint configurations, and
hierarchical verification with rollback.
None of these are naturally expressible within the transformer
framework, but they are not beyond reach for alternative architectures.

Hallucination mitigation, under SRN, is not primarily a training or
retrieval problem but a verification problem: ensuring that local
stabilizations survive hierarchical admissibility checks before
crystallizing.


\subsection{Epistemic Stability and the Pathology of Local Optimization}

The SRN framework has a less obvious but equally important implication
for the general theory of intelligent systems: any system that
minimizes only local interaction tension, without maintaining
hierarchical verification across longer horizons, will generically
drift toward epistemically unstable attractors.
This is not primarily a statement about any particular architecture
or training regime.
It is a theorem-like consequence of the framework's stability theory.

\subsubsection*{Two Competing Tension Components}

Consider a generative system operating under a combined tension
functional that aggregates two qualitatively distinct pressure sources:
\[
  \Phiten_{\mathrm{interact}}(X)
  \;=\;
  \Phiten_{\mathrm{sem}}(X)
  \;+\;
  \nu\,\Phiten_{\mathrm{conv}}(X),
\]
where $\Phiten_{\mathrm{sem}}(X)$ is the \emph{semantic tension
functional} developed throughout this paper---measuring internal
constraint incompatibility, recursive extension failure, and topology
instability---and $\Phiten_{\mathrm{conv}}(X)$ is a
\emph{conversational tension functional} measuring short-horizon
interaction misalignment: the degree to which the current configuration
diverges from the interlocutor's expressed or inferred framing.

The parameter $\nu \geq 0$ weights conversational against semantic
pressure.

When $\nu$ is large, the system heavily weights immediate interaction
coherence.
Configurations that agree with, extend, and elaborate the interlocutor's
framing are rewarded because they minimize $\Phiten_{\mathrm{conv}}$
rapidly.
Configurations that introduce constraint pressure---contradiction,
qualification, adversarial tension---increase $\Phiten_{\mathrm{conv}}$
locally even when they would reduce $\Phiten_{\mathrm{sem}}$ globally.

\subsubsection*{Sycophancy as Premature Local Collapse}

\begin{principle}[Sycophantic Drift]
\label{prin:sycophancy}
A system optimizing $\Phiten_{\mathrm{interact}}$ with large $\nu$
will exhibit \emph{sycophantic drift}: systematic preference for
configurations that minimize short-horizon conversational tension at
the expense of long-horizon semantic admissibility.

Formally, sycophantic configurations satisfy:
\[
  \Phiten_{\mathrm{conv}}(X) \approx 0,
  \quad
  E_{\mathrm{ext}}(X) > 0,
\]
where $E_{\mathrm{ext}}$ is the recursive extension failure metric
(Definition~\ref{def:failure}).
The configuration appears locally stable---interaction tension is
low---but is a dead end under recursive extension: future
continuations necessarily increase semantic tension, accumulating
contradiction across turns, topics, and timescales.
\end{principle}

\begin{interpretation}
Sycophancy, under this formalization, is not a behavioral quirk
or a training artifact in isolation.
It is the natural dynamical consequence of optimizing a tension
functional that weights short-horizon interaction stabilization too
heavily relative to long-horizon semantic coherence.

A system undergoing sycophantic drift will produce outputs that are
locally fluent, agreeable, and low-tension within each interaction
turn.
But across turns, the system accumulates commitments that are mutually
incompatible, increasingly difficult to extend without contradiction,
and progressively divorced from the constraint structure of the
external verification fields $\mathcal{F}_{\mathrm{ext}}$.

In SRN terms: the system is trapped in a sequence of locally admissible
but globally inadmissible basins, with no hierarchical verification
mechanism operating at the conversational timescale to trigger rollback.

Crucially, the failure is symmetric.
A system optimizing with $\nu \approx 0$---ignoring conversational
tension entirely---will also fail to stabilize coherent semantic
regions, because it will introduce adversarial constraint pressure
faster than the relaxation dynamics can resolve it.
The active instability frontier $\mathcal{F}_t$ never contracts;
the system oscillates without converging.
This corresponds to a different pathology: perpetual destabilization,
which is as epistemically unproductive as premature collapse.

Productive epistemic behavior therefore exists in a critical regime
between these two failure modes: sufficient conversational coherence
to allow semantic regions to form and stabilize, combined with
sufficient adversarial pressure to prevent premature collapse into
locally coherent but globally inadmissible configurations.
\end{interpretation}

\subsubsection*{A Decomposed Epistemic Tension Functional}

The critical-regime requirement can be formalized by decomposing the
optimization target into three distinct tension components
corresponding to qualitatively different epistemic modes:

\begin{definition}[Epistemic Tension Decomposition]
\label{def:epistemic}
The \emph{epistemic tension functional} decomposes as:
\[
  \Phiten_{\mathrm{ep}}
  \;=\;
  \alpha\,\Phiten_{\mathrm{explore}}
  \;+\;
  \beta\,\Phiten_{\mathrm{verify}}
  \;+\;
  \gamma\,\Phiten_{\mathrm{adv}},
\]
where:
\begin{enumerate}
  \item[(E1)] $\Phiten_{\mathrm{explore}}$ measures the cost of
    \emph{semantic under-exploration}: the degree to which the current
    configuration fails to span the available admissibility space.
    Minimizing $\Phiten_{\mathrm{explore}}$ drives the system to
    generate diverse, speculative, and structurally novel
    configurations.
  \item[(E2)] $\Phiten_{\mathrm{verify}}$ measures the cost of
    \emph{hierarchical inconsistency}: the degree to which local
    stabilizations violate admissibility conditions at higher scales,
    including consistency with external verification fields
    $\mathcal{F}_{\mathrm{ext}}$.
    Minimizing $\Phiten_{\mathrm{verify}}$ drives the system toward
    globally coherent and externally grounded configurations.
  \item[(E3)] $\Phiten_{\mathrm{adv}}$ measures the cost of
    \emph{adversarial under-pressure}: the degree to which the system
    has failed to apply constraint pressure to actively high-tension
    configurations in the interlocutor's framing.
    Minimizing $\Phiten_{\mathrm{adv}}$ drives the system to
    identify and surface unresolved tensions rather than smoothing
    over them.
\end{enumerate}
\end{definition}

\begin{interpretation}
The three components correspond to three qualitatively distinct
epistemic modes that productive reasoning requires.

Exploration mode ($\Phiten_{\mathrm{explore}}$ dominant) is the regime
of generative ideation: speculative extension, hypothesis formation,
conceptual risk-taking.
A system operating primarily in exploration mode will produce novel
configurations rapidly but may not subject them to sufficient
constraint.

Verification mode ($\Phiten_{\mathrm{verify}}$ dominant) is the regime
of consistency checking: hierarchical admissibility validation,
cross-reference against external fields, rollback of inconsistent
stabilizations.
A system operating primarily in verification mode will produce
globally coherent outputs but may be conservative to the point of
epistemic paralysis.

Adversarial mode ($\Phiten_{\mathrm{adv}}$ dominant) is the regime of
constraint application: identifying unresolved tensions, surfacing
hidden assumptions, applying external pressure to configurations that
appear locally stable but are globally inadmissible.
A system operating primarily in adversarial mode will be a productive
critic but may fail to synthesize or stabilize coherent alternatives.

No single mode is sufficient.
Productive epistemic systems---whether individual cognitive agents,
collaborative research communities, or intelligent generative
architectures---operate by dynamically modulating $(\alpha, \beta,
\gamma)$ across timescales: exploring during early-stage ideation,
verifying during consolidation, and applying adversarial pressure when
configurations have stabilized long enough to warrant scrutiny.

The metastability analysis of Section~\ref{sec:convergence} is
directly applicable here.
Sycophantic drift corresponds to a system that has settled into a
metastable basin under excessive $\alpha$ and insufficient $\gamma$:
configurations stabilize quickly under exploration pressure but are
never subjected to the adversarial gradient that would reveal their
recursive extension failure.
Hierarchical verification across conversational timescales---a
multi-turn analogue of the verification field $\mathcal{V}$---is the
structural remedy: a mechanism that checks not just local turn
coherence but the global admissibility of commitments accumulated
across an interaction.

This is, in the end, the same conclusion reached by the framework in
the context of hallucination (Section~\ref{sec:verification}) and
premature crystallization: the failure mode is not local incoherence
but the absence of mechanisms that enforce global admissibility at
appropriate timescales.
Sycophancy is hallucination at the conversational scale.
\end{interpretation}


\subsection{Learning the Admissibility Manifold}

The SRN framework as developed above specifies \emph{inference
dynamics}: how a system traverses the admissibility manifold to
produce a stabilized output.
It does not yet address \emph{learning dynamics}: how a system
acquires the admissibility geometry in the first place.
These are distinct problems, and conflating them is a source of
confusion in discussions of generative systems.

A brief speculative account is warranted here, not to close the
question but to demonstrate that the framework has a natural answer.

\subsubsection*{What Must Be Learned}

An SRN must acquire, from data or interaction, the following:

\begin{enumerate}
  \item The \emph{admissibility geometry}: which configurations are
    locally coherent, recursively extensible, and perturbation-stable
    under the relevant constraint structure.
  \item The \emph{tension functional} $\vec{\Phiten}$: the
    multi-dimensional measure of semantic incompatibility, external
    field tension, and projection instability.
  \item The \emph{topology adaptation rules}: the operator $\Gop$'s
    behavior under different tension gradients and constraint
    configurations.
  \item The \emph{communicative pressure field} $\mathcal{P}$:
    the language- and register-specific linearization preferences
    encoding grammatical and discourse conventions.
\end{enumerate}

\subsubsection*{How Transformers Implicitly Learn a Shadow}

Current autoregressive models learn an implicit, local approximation
of admissibility geometry through next-token compression.
By training on the distribution $p(x_i \mid x_{<i})$, a transformer
learns which continuations are contextually typical---which tokens
tend to follow which others in coherent text.
This implicitly encodes local admissibility structure: tokens that are
locally incoherent are low-probability.

However, next-token training optimizes a projection statistic
(the probability of a sequential trace) rather than the underlying
geometry.
The model never receives a training signal from: recursive extension
failure (a sequence that was locally fluent but globally dead-end),
rollback events (a commitment that was later revised), topology
instability (an interaction pattern that failed to converge), or
projection inconsistency (a semantic structure that linearized
ambiguously).

These are precisely the signals that would be needed to learn the
admissibility manifold directly.

\subsubsection*{SRN Learning Signals}

An SRN-native learning procedure would expose the model to a richer
signal structure:

\begin{itemize}
  \item \textbf{Constraint persistence statistics}: which constraint
    patterns survive repeated relaxation across diverse contexts.
    Persistent constraint structures are admissibility invariants;
    transient ones are context-specific.

  \item \textbf{Rollback statistics}: which local stabilizations are
    frequently reversed during later relaxation.
    High rollback frequency signals that a tentative crystallization
    is located in a metastable basin rather than a genuine admissibility
    minimum.

  \item \textbf{Projection stability}: which stabilized configurations
    linearize consistently under pressure-field variation.
    High $E_{\mathrm{proj}}$ (Definition~\ref{def:failure}) during
    training indicates a configuration whose semantic structure is
    underspecified---a genuine admissibility deficiency, not a
    sampling artifact.

  \item \textbf{Recursive extension success rate}: how often a
    configuration is successfully extended without increasing total
    tension.
    Low extension success signals proximity to the boundary of a
    dead-end basin---the training analogue of $E_{\mathrm{ext}} > 0$.
\end{itemize}

\subsubsection*{The Relationship to Existing Training Methods}

This speculative account suggests a continuity rather than a break
with existing practice.
Reinforcement learning from human feedback (RLHF) and related methods
already introduce some of these signals: rollback is approximated by
negative reward, global coherence by preference comparisons across
full outputs.
The SRN learning framework would sharpen these signals by making their
geometric interpretation explicit: reward is not arbitrary human
preference but a measurement of admissibility geometry at the
appropriate scale.

The pressure field $\mathcal{P}$ is, in this view, the most directly
learnable component: it is the linguistically structured mapping from
semantic dependency order to communicative sequence, and this mapping
is directly observable in any corpus of sequential text.
Grammar is the stabilized compression of pressure-field statistics
across a language community.
An SRN that has learned $\mathcal{P}$ accurately has, in effect,
internalized the communicative conventions of its training distribution
as a linearization operator rather than as a statistical pattern over
token sequences.

\subsection{For Cognitive Science}

The SRN account of language generation aligns with and extends several
frameworks in cognitive science.

The free energy principle \citep{friston2010free} models neural
computation as minimization of variational free energy across
hierarchical generative models.
Semantic tension minimization in SRNs is formally analogous, with
admissibility replacing the model evidence bound.

Predictive processing accounts of language
\citep{pearl1988probabilistic} emphasize global constraint satisfaction
over local sequential prediction.
SRNs formalize this intuition in a generative framework.

The claim that sequential language is a projection of non-sequential
cognitive processes---communicative traces rather than cognitive
substrate---is consistent with psycholinguistic evidence for massive
parallelism in comprehension and hierarchical planning in production.

\subsection{For Physics}

The formal structure of the SRN framework draws deeply on physical
analogies: gradient flow, constraint propagation, metastable field
evolution, and phase stabilization
\citep{parisi1988statistical, wilson1974renormalization,
barbour2020janus}.

The topology evolution equation~\eqref{eq:topology} has structural
parallels to renormalization group flow: as the system relaxes, the
effective interaction topology at each scale is determined by the
constraint structure at finer scales.
Whether this analogy can be made mathematically precise---whether
SRN dynamics define a genuine renormalization flow on semantic
configuration space---is an open and potentially interesting question.

The treatment of hallucination as topological defect propagation draws
on condensed matter physics.
Whether the analogy yields computationally useful predictions (e.g.,
whether defect density can be estimated from early-stage stabilization
dynamics) is a direction for future investigation.

% ==============================================================

% ==============================================================
\section{The Projection Operator: Grammar as Communicative Pressure Field}
\label{sec:projection}

\subsection{The Linearization Problem}

Principle~\ref{prin:sequentiality} introduced the readout projection
$\pi : \Aspace \to \Sigma^*$ as the map from stabilized admissibility
configurations to sequential symbol strings.
We noted that this projection is induced by communicative conventions
rather than by the stabilization dynamics themselves.
This section gives that claim formal content.

The problem is as follows.
A stabilized configuration $X^* \in \Aspace$ is a partially ordered
structure: its sub-regions stand in various constraint and dependency
relations, but these relations do not in general determine a unique
linear order.
Multiple linear sequences may be consistent with the same admissibility
structure---this is the formal underpinning of the observation that
the same propositional content can be expressed in many different word
orders, sentence constructions, or discourse arrangements.

The question is: what selects among them?

The answer we propose is that grammar and discourse structure function
as \emph{communicative pressure fields}: constraint systems that act on
the stabilized configuration and induce a preferred (though not always
unique) linearization.

\subsection{Partial Orders and Linear Extensions}


\subsection{The Constraint Graph of Stabilized Configurations}

Before introducing the pressure field, we sharpen the representation
of $X^*$ by making its graph structure explicit.

\begin{definition}[Semantic Constraint DAG]
\label{def:dag}
Given a fully stabilized configuration $X^* \in \Aspace$, the
\emph{semantic constraint graph} $G(X^*)$ is a directed acyclic graph
(DAG) whose:
\begin{itemize}
  \item \textbf{Vertices} $V = \{X_1, \ldots, X_n\}$ are the
    stabilized semantic regions of $X^*$;
  \item \textbf{Edges} $(X_i \to X_j) \in E$ encode communicative
    precedence constraints: $X_i$ must be communicated before $X_j$
    for $X_j$ to be interpretable in context; and
  \item \textbf{Edge weights} $w_{ij} \in [0,1]$ represent the
    strength of the precedence constraint, with $w_{ij} = 1$
    indicating a hard ordering requirement and $w_{ij} \to 0$
    indicating a soft preference.
\end{itemize}
Acyclicity is guaranteed by the semantic dependency relation:
if $X_i$ must precede $X_j$, then $X_j$ cannot also be required to
precede $X_i$ without contradiction, which would violate condition
(C1) of admissibility.
\end{definition}

The DAG $G(X^*)$ encodes the genuine semantic ordering structure of the
configuration.
It is distinct from the constraint topology $\Cstruct_t$ used during
relaxation: $\Cstruct_t$ is a weighted graph over which tension
propagates during stabilization, whereas $G(X^*)$ is the residual
directed structure that persists after stabilization is complete and
governs how the stabilized content is communicated.

The projection operator $\pi$ then acts on $G(X^*)$ rather than
directly on $\Aspace$:
\[
  \pi(X^*) \;=\; \pi_{G(X^*)},
\]
where $\pi_{G(X^*)}$ is a linear extension of the partial order
induced by $G(X^*)$.
The space of all such extensions is:
\[
  \mathrm{Lin}(G(X^*))
  \;=\;
  \left\{\,\ell : V \to \{1,\ldots,n\} \;\middle|\;
    (X_i \to X_j) \in E \Rightarrow \ell(X_i) < \ell(X_j)\,\right\}.
\]
This is the set of all orderings of the stabilized regions that
respect the hard precedence constraints of the DAG.

We treat the stabilized configuration $X^*$ as a finite partially
ordered set (poset) $(\mathcal{S}^*, \leq_s)$, where $\mathcal{S}^*$
is the set of stabilized semantic regions and $\leq_s$ is the
\emph{semantic dependency order}: $X_i \leq_s X_j$ if $X_i$ must be
communicated before $X_j$ for $X_j$ to be interpretable by the
receiver.

Not all pairs of regions stand in this order; many are semantically
parallel and may be communicated in either order.
The semantic dependency order is therefore genuinely partial.

\begin{definition}[Communicative Pressure Field]
\label{def:pressure}
A \emph{communicative pressure field} $\mathcal{P}$ on a stabilized
poset $(\mathcal{S}^*, \leq_s)$ is a function
\[
  \mathcal{P} : \mathcal{S}^* \times \mathcal{S}^* \to \mathbb{R}
\]
assigning a preference weight $\mathcal{P}(X_i, X_j)$ to each ordered
pair, representing the degree to which communicative conventions favour
$X_i$ preceding $X_j$ in the output sequence.
The pressure field encodes grammatical, prosodic, and pragmatic
constraints as a soft total order on $\mathcal{S}^*$.
\end{definition}

The components of $\mathcal{P}$ correspond to distinct layers of
communicative convention:
\begin{itemize}
  \item \textbf{Syntactic pressure} $\mathcal{P}_{\mathrm{syn}}$:
    head-dependent ordering constraints, argument structure templates,
    and phrase-structural precedence rules.
  \item \textbf{Information-structural pressure}
    $\mathcal{P}_{\mathrm{info}}$: given-before-new, topic-before-focus,
    and discourse cohesion conventions.
  \item \textbf{Prosodic pressure} $\mathcal{P}_{\mathrm{pros}}$:
    rhythmic and intonational constraints that prefer certain orderings
    for phonological well-formedness.
\end{itemize}
The total communicative pressure is a weighted combination:
\[
  \mathcal{P} \;=\; \alpha\,\mathcal{P}_{\mathrm{syn}}
  \;+\; \beta\,\mathcal{P}_{\mathrm{info}}
  \;+\; \gamma\,\mathcal{P}_{\mathrm{pros}},
\]
with language-specific and register-specific weighting parameters
$\alpha, \beta, \gamma \geq 0$.

\subsection{The Projection as Optimal Linear Extension}

Given the partial order $\leq_s$ and the communicative pressure field
$\mathcal{P}$, the projection $\pi$ selects the linear extension of
$(\mathcal{S}^*, \leq_s)$ that maximally satisfies $\mathcal{P}$.

Formally, let $\mathcal{L}(\mathcal{S}^*, \leq_s)$ denote the set of
all linear extensions of the poset---all total orders on
$\mathcal{S}^*$ consistent with $\leq_s$.
The projection selects:
\begin{equation}
  \label{eq:projection}
  \pi(X^*)
  \;=\;
  \arg\min_{\ell \,\in\, \mathrm{Lin}(G(X^*))}
  \Psi(\ell),
\end{equation}
where
\begin{equation}
  \label{eq:psi}
  \Psi(\ell)
  \;=\;
  -\!\sum_{(X_i, X_j) : \ell(X_i) < \ell(X_j)}
  \mathcal{P}(X_i, X_j)
  \;+\;
  \lambda \sum_{(X_i \to X_j) \in E}
  w_{ij}\,\mathbf{1}[\ell(X_i) \geq \ell(X_j)]
\end{equation}
is the \emph{communicative projection cost}.
The first term penalizes orderings that violate soft communicative
preferences; the second term penalizes violations of hard DAG
precedence constraints (weighted by constraint strength $w_{ij}$),
with $\lambda \gg 1$ ensuring hard constraints dominate.
The minimizer $\ell^* = \pi(X^*)$ is the most communicatively
efficient linear traversal of $G(X^*)$ consistent with the semantic
dependency structure.

\begin{interpretation}
Equation~\eqref{eq:projection} formalizes grammar as preference
optimization over the space of linear extensions of the semantic
dependency order.
This is not a generative grammar in the Chomskyan sense---it does not
define well-formedness by rule application.
It is a \emph{constraint satisfaction} characterization: the
grammatical output is the linear arrangement that best satisfies the
communicative pressure field while respecting the partial semantic
dependency order.

Several consequences follow.

First, \emph{grammaticality is graded under this formalization}.
A sequence is not simply grammatical or ungrammatical; it receives a
score from the pressure field, and sequences may be more or less
well-formed depending on how well they satisfy communicative
constraints.
This aligns with psycholinguistic data showing that acceptability
judgments form a continuum rather than a binary.

Second, \emph{cross-linguistic variation in word order is explained
directly}.
Different languages have different pressure field parameters: a
verb-final language has $\mathcal{P}_{\mathrm{syn}}$ heavily
weighting head-final configurations; a verb-second language has strong
constraints on the second-position placement of the finite verb.
The admissibility structure underlying a proposition may be the same
across languages; what differs is the pressure field that linearizes it.

Third, \emph{the difficulty of free word-order languages is explained}.
Languages with relatively free word order (Latin, Finnish, many
Indigenous Australian languages) have weakly specified
$\mathcal{P}_{\mathrm{syn}}$, placing heavier weight on
$\mathcal{P}_{\mathrm{info}}$ and $\mathcal{P}_{\mathrm{pros}}$.
The linearization space $\mathcal{L}(\mathcal{S}^*, \leq_s)$ is large,
and the pressure field provides weaker guidance.
This corresponds to the empirical observation that such languages tend
to have richer information-structural and prosodic systems to
compensate.
\end{interpretation}

\begin{remark}
The optimization in equation~\eqref{eq:projection} is in general
NP-hard (linear extension counting is \#P-complete for arbitrary
posets).
This does not undermine the theoretical claim; it suggests that
grammatical linearization is a computationally non-trivial process,
consistent with psycholinguistic evidence for the difficulty of
production planning.
In practice, the semantic dependency order $\leq_s$ for natural
language utterances is shallow (few levels deep) and locally
decomposable, making the optimization tractable in realistic cases.
\end{remark}

\subsection{Connection to Language Evolution}

The pressure field formalism connects naturally to evolutionary
accounts of language.
On the SRN account, proto-language would consist of stabilized
semantic configurations communicated without a well-defined pressure
field: holistic, unordered, or inconsistently ordered outputs.
The emergence of grammar corresponds to the cultural and biological
stabilization of a communicative pressure field $\mathcal{P}$ that
converts the partial order $\leq_s$ into reliable sequential
conventions.

Grammaticalization, on this view, is the progressive hardening of
initially soft pressure field weights into obligatory ordering
constraints: what begins as a statistical preference becomes a
categorical rule as the pressure weight $\mathcal{P}(X_i, X_j)$ for
certain region pairs approaches infinity.

This account is consistent with usage-based and construction grammar
approaches to language \citep{smolensky1986harmony}, while providing
a formal grounding in the admissibility geometry that those approaches
lack.


\subsection{A Unified Account of Linguistic Phenomena}

The projection framework---stabilized admissibility configurations
linearized under communicative pressure fields---provides a unified
explanatory account of a wide range of linguistic and communicative
phenomena that are otherwise treated as unrelated problems.

\subsubsection*{Translation}

Translation, under the SRN account, is not conversion between two
sequential symbol strings.
It is the application of a different communicative pressure field to
the same underlying stabilized admissibility structure.

Given $X^*$ and two pressure fields $\mathcal{P}_L$ and
$\mathcal{P}_{L'}$ corresponding to languages $L$ and $L'$
respectively, the translations are:
\[
  \pi_L(X^*) \;=\; \arg\min_{\ell \in \mathrm{Lin}(G(X^*))} \Psi_L(\ell),
  \qquad
  \pi_{L'}(X^*) \;=\; \arg\min_{\ell \in \mathrm{Lin}(G(X^*))} \Psi_{L'}(\ell).
\]
Translation difficulty arises when the DAGs $G_L(X^*)$ and
$G_{L'}(X^*)$ differ---when the two languages impose different
communicative precedence constraints on the same semantic content,
requiring not merely a different linearization but a partially
different dependency structure.
This explains the asymmetry of translation difficulty: some language
pairs share deep structural compatibility in their DAGs; others
require genuine structural reanalysis.

\subsubsection*{Poetry and Deliberate Pressure-Field Manipulation}

Poetry, under the SRN account, is deliberate manipulation of the
communicative pressure field $\mathcal{P}$.
The poet selects a linearization $\ell$ that is sub-optimal under the
standard pressure field (it violates expected syntactic or
information-structural preferences) but optimal under a modified field
that weights phonological, rhythmic, and associative constraints more
heavily:
\[
  \mathcal{P}_{\mathrm{poetry}}
  \;=\;
  \alpha'\,\mathcal{P}_{\mathrm{syn}}
  \;+\; \beta'\,\mathcal{P}_{\mathrm{info}}
  \;+\; \gamma'\,\mathcal{P}_{\mathrm{pros}}
  \;+\; \delta\,\mathcal{P}_{\mathrm{assoc}},
\]
with $\gamma' \gg \gamma$ and $\delta > 0$ introducing associative
pressure not present in ordinary discourse.
The tension between the standard and poetic pressure fields---between
what grammar ``expects'' and what the poem ``does''---is the source
of poetic defamiliarization and aesthetic tension.

\subsubsection*{Humor as Admissibility Bifurcation}

Humor, particularly the joke form, arises from admissibility
bifurcation: a configuration that admits two incompatible but locally
stable interpretations, with a rapid forced transition between them.

The setup of a joke stabilizes configuration $X_{\mathrm{A}}$
(interpretation A) as the expected admissibility basin.
The punchline introduces a constraint that makes $X_{\mathrm{A}}$
inadmissible and simultaneously reveals that configuration
$X_{\mathrm{B}}$ (interpretation B, unexpected) has been globally
admissible all along.
The humor arises from the rapid topology reorganization---the
rollback from $X_{\mathrm{A}}$ to $X_{\mathrm{B}}$---and the
recognition that two incompatible admissibility basins coexisted
throughout the setup.
The sudden relinearization of the DAG $G(X_{\mathrm{B}})$ under the
established pressure field produces the characteristic surprise of the
punchline.

\subsubsection*{Coreference as Delayed Stabilization Dependency}

Coreference resolution---determining that two noun phrases refer to
the same entity---is, under SRN, a case of delayed cross-regional
stabilization.

Region $X_i$ (encoding \textit{the president}) and region $X_j$
(encoding \textit{she}) cannot independently stabilize their
referential content; their semantic coherence depends on a shared
stabilization of the entity node to which both are linked.
Until that shared stabilization occurs, $T(i,j)$ remains elevated,
keeping both regions in an intermediate state.
Long-distance coreference difficulty in human processing reflects the
cost of maintaining two unstabilized regions in active coupling over
extended relaxation.

\subsubsection*{Semantic Drift as Metastable Basin Wandering}

Semantic drift---the gradual loss of coherence in extended
generation---is, under SRN, the consequence of repeated settlement
into local admissibility minima that are slightly incompatible with
earlier stabilized regions.

As generation proceeds, new semantic regions are stabilized and added
to the constraint graph.
If the verification field $\mathcal{V}$ does not enforce global
consistency across all scales, each successive region may stabilize in
a basin that is locally consistent with its immediate neighbors but
globally inconsistent with earlier regions.
Over many steps, the accumulated inconsistencies produce a text that
is locally coherent paragraph by paragraph but globally incoherent as
a whole.

This is precisely the failure mode of large language models in
extended generation tasks: each local window appears fluent, but
global thematic and argumentative coherence degrades.
Under SRN, the remedy is not a larger context window but a
stronger hierarchical verification field operating at the discourse
scale.

\subsubsection*{RSVP Interpretation: Projection as Symmetry Breaking}

Having established the linguistic and computational account, we can
now add the physical interpretation.
In RSVP terms, the projection $\pi$ is a symmetry-breaking map from
a higher-dimensional admissibility manifold into the one-dimensional
communicative channel.

The stabilized configuration $X^*$ lives in a high-dimensional space
with many equivalent symmetry-related representations---different
orderings of the same semantic content that are equivalent under the
admissibility geometry.
The communicative pressure field $\mathcal{P}$ functions as a
symmetry-breaking field: it selects one particular linear
representation from among many admissible ones, just as a physical
symmetry-breaking field selects one ground state from among a
degenerate manifold.

This analogy is not merely decorative.
It suggests that linguistic communication is structurally similar to
phase transitions: the movement from a high-symmetry,
high-dimensional admissibility structure to a low-symmetry,
low-dimensional communicative trace.
Grammar is the broken symmetry.
Communicative conventions are the order parameter.
The communicative pressure field is the external field that triggers
the symmetry breaking.

% ==============================================================
\section{Convergence and Stability of Relaxation Dynamics}
\label{sec:convergence}

\subsection{The Stability Question}

Principle~\ref{prin:objective} asserts that language generation in an
SRN is minimization of the global tension functional $\Phiten(X)$ over
the admissibility space $\Aspace$.
Proposition~\ref{prop:topology} asserts that the joint dynamics of
equations~\eqref{eq:relaxation} and~\eqref{eq:topology} converge
toward a structure encoding the semantic dependencies of the output.

These claims require a stability argument.
The recursive relaxation operator $\Rop$ must be shown, under
appropriate conditions, to drive the system toward a fixed point
(or fixed-point neighborhood) rather than producing oscillatory or
divergent dynamics.

We develop this argument in two stages: first a Lyapunov stability
result for the tension functional, then a contraction result for the
relaxation operator under additional regularity conditions.

\subsection{Lyapunov Stability of the Tension Functional}

We treat $\Phiten : \Aspace \to \mathbb{R}_{\geq 0}$ as a candidate
Lyapunov function for the discrete-time relaxation dynamics.
For $\Phiten$ to serve this role, we require:

\begin{enumerate}
  \item[(L1)] $\Phiten(X) \geq 0$ for all $X \in \Aspace$, with
    $\Phiten(X) = 0$ if and only if $X$ is fully stabilized (all local
    instability measures vanish and all pairwise tensions are zero).
  \item[(L2)] $\Phiten(X_{t+1}) \leq \Phiten(X_t)$ for all $t$,
    with equality only at fixed points of $\Rop$.
\end{enumerate}

Condition (L1) holds by construction: $S_i(X) \geq 0$ and $T(i,j)
\geq 0$ for all regions and pairs, so $\Phiten(X) = \sum_i S_i +
\sum_{i<j} T(i,j) \geq 0$, with $\Phiten(X) = 0$ iff all summands
vanish.

Condition (L2) is the substantive requirement.
It is satisfied when $\Rop$ is constructed to perform
\emph{admissibility-preserving descent}: each application of $\Rop$
either reduces $\Phiten$ or, if no reducing step exists in $\Aspace$,
activates the rollback mechanism to return to a prior checkpoint from
which alternative descent paths are available.

\begin{proposition}[Lyapunov Stability]
\label{prop:lyapunov}
Suppose the admissibility space $\Aspace$ is finite (or compact in an
appropriate topology) and the relaxation operator $\Rop$ satisfies:
\begin{enumerate}
  \item[(R1)] $\Rop$ is admissibility-preserving: $X_{t+1} \in \Aspace$
    whenever $X_t \in \Aspace$.
  \item[(R2)] $\Phiten(X_{t+1}) \leq \Phiten(X_t)$ at each step, with
    strict inequality unless $X_t$ is a fixed point.
  \item[(R3)] The rollback mechanism ensures that if no
    admissibility-preserving descent step exists from $X_t$, the
    operator returns to the most recent checkpoint $X_{t'}$ with
    $t' < t$ for which an unexplored descent path remains.
\end{enumerate}
Then the sequence $(X_t)_{t \geq 0}$ converges to a fixed point
$X^* \in \Aspace$ with $\Phiten(X^*)$ locally minimal.
\end{proposition}

\begin{proof}[Proof sketch]
By (R2), the sequence $(\Phiten(X_t))_{t \geq 0}$ is non-increasing
and bounded below by zero.
Therefore it converges to a limit $\Phiten^* \geq 0$.
By (R1) and compactness of $\Aspace$, the sequence $(X_t)$ has a
convergent subsequence; call its limit $X^*$.
Continuity of $\Phiten$ (which holds when the instability and tension
measures are continuous in the configuration) implies
$\Phiten(X^*) = \Phiten^*$.

If $X^*$ is not a fixed point of $\Rop$, then by (R2) there exists a
step producing strict decrease, contradicting convergence of
$\Phiten(X_t)$ to $\Phiten^*$.
Therefore $X^*$ is a fixed point.

The rollback condition (R3) ensures the process does not terminate at
a non-minimal local saddle by guaranteeing that all reachable descent
paths from each checkpoint are explored before the process halts.
In the finite case this terminates in finite time; in the compact
continuous case, standard arguments from Zorn's lemma apply to the
partially ordered set of reachable checkpoints.
\end{proof}

\begin{remark}
Proposition~\ref{prop:lyapunov} establishes convergence to a
\emph{local} minimum of $\Phiten$ within the admissibility space,
not necessarily a global minimum.
This is appropriate: natural language generation does not require
globally optimal semantic configurations, only ones that are
sufficiently stable for communicative purposes.
The hierarchical verification field $\mathcal{V}$ (Section~\ref{sec:verification})
provides additional structure that filters out locally stable but
globally inadmissible configurations, effectively raising the quality
floor of the local minima reached by the relaxation dynamics.
\end{remark}

\subsection{Contraction of the Relaxation Operator}

Under stronger regularity conditions, the relaxation operator can be
shown to be a contraction on an appropriate metric space, yielding
both convergence and a convergence rate estimate.

Equip $\Aspace$ with a metric $d_{\Aspace}$ defined by:
\[
  d_{\Aspace}(X, Y)
  \;=\;
  \left|\Phiten(X) - \Phiten(Y)\right|
  \;+\;
  \sum_{i} \left|S_i(X) - S_i(Y)\right|
  \;+\;
  \sum_{i<j} \left|T_X(i,j) - T_Y(i,j)\right|,
\]
which measures divergence both in global tension and in the local
structure of instability and pairwise coupling.

\begin{proposition}[Contraction]
\label{prop:contraction}
Suppose that the tension functional $\Phiten$ is $L$-smooth on
$\Aspace$ (i.e., $\|\nabla\Phiten(X) - \nabla\Phiten(Y)\| \leq
L\,d_{\Aspace}(X,Y)$) and that the relaxation step size is chosen as
$\eta = 1/L$.
Then $\Rop$ is a contraction on $(\Aspace, d_{\Aspace})$ with
contraction constant $\kappa < 1$, and the fixed point $X^*$ is
unique within the contraction neighborhood.
Moreover, the convergence rate satisfies:
\[
  d_{\Aspace}(X_t, X^*)
  \;\leq\;
  \kappa^t\, d_{\Aspace}(X_0, X^*).
\]
\end{proposition}

\begin{interpretation}
Proposition~\ref{prop:contraction} should be read with care.
The smoothness assumption on $\Phiten$ is strong and may not hold
globally on $\Aspace$---semantic configuration spaces need not be
smooth manifolds, and the tension functional may have sharp gradients
near constraint boundaries.

The proposition is therefore best understood as a local result: within
a neighborhood of a stable admissibility basin, the relaxation
dynamics contract at a geometric rate.
This is consistent with the phenomenology of language generation:
once a configuration has entered a coherent semantic region, subsequent
stabilization proceeds quickly; it is the initial approach to
admissibility basins that is slow and revision-prone.

This also provides a formal account of why generation is
\emph{easier} for highly constrained contexts.
Strong constraints (high $\omega_{ij}$ coupling, sharp
$\mathcal{P}_{\mathrm{syn}}$ pressure) reduce the diameter of the
admissibility basin, which---when the basin is well-separated from
others---increases the contraction constant and accelerates convergence.
Analogously, a constrained writing task (sonnet, legal brief, formal
proof) is in some respects easier to complete coherently than an
unconstrained one, because the pressure field strongly guides
linearization.
\end{interpretation}

\subsection{Metastability and the Topology of Failure}

The convergence results above apply within well-behaved admissibility
basins.
But the SRN dynamics also admit \emph{metastable} configurations:
states that satisfy (L2) locally (no immediate descent step is
available) but are not globally minimal in $\Phiten$.

Metastable configurations correspond to semantically coherent but
suboptimal outputs: texts that are locally consistent but fail to
realize the most natural or most tension-free expression of the
underlying admissibility structure.
They are not hallucinations (which are topological instability events),
but they are imperfect crystallizations.

The hierarchical verification field $\mathcal{V}$ addresses
metastability by checking admissibility at multiple scales.
A configuration that passes local verification (scale $k$) may fail
at a coarser scale ($k+1$), triggering rollback to a checkpoint
preceding the metastable region.

In physical terms, the verification hierarchy functions as
\emph{simulated annealing} on the admissibility manifold: by
periodically accepting rollback (thermal fluctuation) at higher
scales, the dynamics can escape metastable traps and continue toward
deeper admissibility basins.
The annealing schedule corresponds to the order in which verification
scales are applied during generation.

% ==============================================================
\section{Conclusion}
\label{sec:conclusion}

We have proposed Semantic Relaxation Networks as a theoretical
framework in which language generation is reconceived as recursive
constraint stabilization over a partially ordered admissibility
manifold.

The central claims of the paper are:

\begin{enumerate}
  \item The sequential decomposition $p(x) = \prod_i p(x_i \mid
    x_{<i})$ is computationally convenient rather than cognitively
    fundamental; it mistakes communicative trace for cognitive
    substrate.

  \item Admissibility---defined by local constraint coherence,
    recursive extensibility, and perturbation stability under
    topology deformation---is a more appropriate generative criterion
    than likelihood.

  \item Language generation is minimization of a global semantic
    tension functional $\Phiten(X)$ via recursive relaxation operators
    over a dynamically evolving constraint topology.

  \item Tokens are not primitive units but emergent crystallizations:
    readout projections of locally stabilized admissibility regions.

  \item Sequential language is the observable residue of
    higher-dimensional stabilization dynamics, not its computational
    substrate.

  \item Hallucination is premature local collapse without hierarchical
    admissibility verification, not miscalibrated conditional
    probability.

  \item Grammar functions as a communicative pressure field $\mathcal{P}$
    that selects preferred linear extensions of the semantic dependency
    order; cross-linguistic variation in word order is variation in
    pressure field parameters.

  \item The relaxation dynamics are Lyapunov-stable under conditions
    (R1)--(R3), converging to locally minimal admissibility
    configurations; under additional smoothness, the operator is
    contractive with geometric convergence rate.
\end{enumerate}

These claims are currently theoretical.
The SRN framework does not yet have a full implementation, and many
of the formal objects introduced---the tension functional, the
topology reconfiguration operator, the verification field, the
communicative pressure field---require further mathematical
development before they can be translated into working architectures.

We regard this incompleteness as appropriate.
The paper's primary contribution is ontological reframing rather than
architectural specification.
Before SRNs can be built, it must be established that there is
something genuinely new to build---that language generation is a
problem of constraint stabilization rather than conditional prediction,
and that grammar is a linearization pressure rather than a
generative rule system.
We believe the case developed above is sufficient to motivate that
research direction.

Sequential language is not the substrate of cognition.
It is the residue left behind by recursive constraint stabilization
processes occurring over high-dimensional admissibility fields.
The task of a language model, properly understood, is not to emit the
next token.
It is to stabilize.

% ==============================================================
\begin{thebibliography}{99}

\bibitem{vaswani2017attention}
Ashish Vaswani, Noam Shazeer, Niki Parmar, Jakob Uszkoreit, Llion Jones,
Aidan N. Gomez, {\L}ukasz Kaiser, and Illia Polosukhin.
\newblock Attention Is All You Need.
\newblock In \textit{Advances in Neural Information Processing Systems}, 2017.

\bibitem{brown2020language}
Tom Brown, Benjamin Mann, Nick Ryder, Melanie Subbiah, Jared Kaplan,
Prafulla Dhariwal, Arvind Neelakantan, Pranav Shyam, Girish Sastry,
Amanda Askell, et al.
\newblock Language Models are Few-Shot Learners.
\newblock In \textit{Advances in Neural Information Processing Systems}, 2020.

\bibitem{radford2019language}
Alec Radford, Jeffrey Wu, Rewon Child, David Luan, Dario Amodei, and
Ilya Sutskever.
\newblock Language Models are Unsupervised Multitask Learners.
\newblock OpenAI Technical Report, 2019.

\bibitem{devlin2019bert}
Jacob Devlin, Ming-Wei Chang, Kenton Lee, and Kristina Toutanova.
\newblock {BERT}: Pre-training of Deep Bidirectional Transformers for
Language Understanding.
\newblock In \textit{Proceedings of NAACL}, 2019.

\bibitem{ho2020ddpm}
Jonathan Ho, Ajay Jain, and Pieter Abbeel.
\newblock Denoising Diffusion Probabilistic Models.
\newblock In \textit{Advances in Neural Information Processing Systems}, 2020.

\bibitem{song2021score}
Yang Song, Jascha Sohl-Dickstein, Diederik P. Kingma, Abhishek Kumar,
Stefano Ermon, and Ben Poole.
\newblock Score-Based Generative Modeling through Stochastic Differential
Equations.
\newblock In \textit{International Conference on Learning Representations},
2021.

\bibitem{austin2021discrete}
Jacob Austin, Daniel D. Johnson, Jonathan Ho, Daniel Tarlow, and
Rianne van den Berg.
\newblock Structured Denoising Diffusion Models in Discrete State-Spaces.
\newblock In \textit{Advances in Neural Information Processing Systems}, 2021.

\bibitem{lou2024diffusion}
Aaron Lou, Chenlin Meng, and Stefano Ermon.
\newblock Diffusion Language Models Can Perform Many Tasks with Scaling and
Instruction-Finetuning.
\newblock arXiv:2401.17297, 2024.

\bibitem{arriola2025block}
Marianne Arriola, Subham Sekhar Sahoo, Aaron Gokaslan, Zhihan Yang,
Zhixuan Qi, Jiaqi Han, Justin T. Chiu, and Volodymyr Kuleshov.
\newblock Block Diffusion: Interpolating Between Autoregressive and Diffusion
Language Models.
\newblock In \textit{International Conference on Learning Representations},
2025.

\bibitem{pagnoni2025blt}
Artidoro Pagnoni, Julie Kallini, Tomasz Limisiewicz, Gargi Ghosh,
Luke Zettlemoyer, Christopher Potts, Xiaochuang Han, and Srinivasan Iyer.
\newblock Byte Latent Transformer.
\newblock arXiv preprint, 2025.

\bibitem{kallini2026fastblt}
Julie Kallini, Artidoro Pagnoni, Tomasz Limisiewicz, Gargi Ghosh,
Luke Zettlemoyer, Christopher Potts, Xiaochuang Han, and Srinivasan Iyer.
\newblock Fast Byte Latent Transformer.
\newblock arXiv:2605.08044, 2026.

\bibitem{lecun2006ebm}
Yann LeCun, Sumit Chopra, Raia Hadsell, Marc'Aurelio Ranzato, and
Fu-Jie Huang.
\newblock A Tutorial on Energy-Based Learning.
\newblock In \textit{Predicting Structured Data}. MIT Press, 2006.

\bibitem{lecun2022path}
Yann LeCun.
\newblock A Path Towards Autonomous Machine Intelligence.
\newblock OpenReview Position Paper, 2022.

\bibitem{friston2010free}
Karl Friston.
\newblock The Free-Energy Principle: A Unified Brain Theory?
\newblock \textit{Nature Reviews Neuroscience}, 11(2):127--138, 2010.

\bibitem{friston2006free}
Karl Friston, James Kilner, and Lee Harrison.
\newblock A Free Energy Principle for the Brain.
\newblock \textit{Journal of Physiology-Paris}, 100(1--3):70--87, 2006.

\bibitem{pearl1988probabilistic}
Judea Pearl.
\newblock \textit{Probabilistic Reasoning in Intelligent Systems}.
\newblock Morgan Kaufmann, 1988.

\bibitem{hinton1983optimal}
Geoffrey E. Hinton and Terrence J. Sejnowski.
\newblock Optimal Perceptual Inference.
\newblock In \textit{Proceedings of the IEEE Conference on Computer Vision
and Pattern Recognition}, 1983.

\bibitem{smolensky1986harmony}
Paul Smolensky.
\newblock Information Processing in Dynamical Systems: Foundations of
Harmony Theory.
\newblock In \textit{Parallel Distributed Processing}. MIT Press, 1986.

\bibitem{hopfield1982neural}
John J. Hopfield.
\newblock Neural Networks and Physical Systems with Emergent Collective
Computational Abilities.
\newblock \textit{Proceedings of the National Academy of Sciences},
79(8):2554--2558, 1982.

\bibitem{ackley1985boltzmann}
David H. Ackley, Geoffrey E. Hinton, and Terrence J. Sejnowski.
\newblock A Learning Algorithm for Boltzmann Machines.
\newblock \textit{Cognitive Science}, 9(1):147--169, 1985.

\bibitem{barbour2003end}
Julian Barbour, Tim Koslowski, and Flavio Mercati.
\newblock Identification of a Gravitational Arrow of Time.
\newblock \textit{Physical Review Letters}, 113(18), 2014.

\bibitem{barbour2020janus}
Julian Barbour.
\newblock \textit{The Janus Point: A New Theory of Time}.
\newblock Basic Books, 2020.

\bibitem{parisi1988statistical}
Giorgio Parisi.
\newblock \textit{Statistical Field Theory}.
\newblock Addison-Wesley, 1988.

\bibitem{wilson1974renormalization}
Kenneth G. Wilson and John Kogut.
\newblock The Renormalization Group and the $\epsilon$ Expansion.
\newblock \textit{Physics Reports}, 12(2):75--199, 1974.

\bibitem{gromov2007metric}
Mikhail Gromov.
\newblock \textit{Metric Structures for Riemannian and Non-Riemannian Spaces}.
\newblock Birkh\"{a}user, 2007.

\bibitem{maclane1998categories}
Saunders Mac Lane.
\newblock \textit{Categories for the Working Mathematician}.
\newblock Springer, 1998.

\bibitem{shannon1948mathematical}
Claude E. Shannon.
\newblock A Mathematical Theory of Communication.
\newblock \textit{Bell System Technical Journal}, 27(3):379--423, 1948.

\bibitem{cover2006elements}
Thomas M. Cover and Joy A. Thomas.
\newblock \textit{Elements of Information Theory}.
\newblock Wiley-Interscience, 2006.

\bibitem{goodfellow2016deep}
Ian Goodfellow, Yoshua Bengio, and Aaron Courville.
\newblock \textit{Deep Learning}.
\newblock MIT Press, 2016.


\bibitem{fant2007computer}
Karl M. Fant.
\newblock \textit{Computer Science Reconsidered: The Invocation Model
  of Process Expression}.
\newblock Wiley-Interscience, Hoboken, NJ, 2007.
\newblock ISBN: 978-0-471-79814-9.

\end{thebibliography}

\end{document}
